\chapter{Concrete examples}\label{parApplications}%
\addcontentsline{itc}{chapter}{\protect\numberline{\thechapter}Concrete examples}

It is now time to see what the results of Chapters~\ref{parTensorization}
and~\ref{parOtherApplicationsOfTensorizationTechniques}
yield for concrete models of statistical physics.
I will try to give rather different types of examples,
so as to illustrate the advantages of working with Hilbertian correlations:
this frame is indeed quite general, as it requires little structure on the models considered.

In \S~\ref{parBackToIsingsModel} we will look back at Ising's model,
seeing how tensorization of Hilbertian decorrelations
improves the results of \S~\ref{parSomeResultsOnIsingsModel},
and what other results are given by the theorems of
\S~\ref{parOtherApplicationsOfTensorizationTechniques}.
We will also consider two kinds of generalizations, namely
when the range of interactions becomes infinite
and when the strength of the interactions is random (\emph{spin glasses}).
In the two next sections we will look at models with continuous states spaces:
first a quite general class of linear models [\S~\ref{parQuadraticModels}],
then a family of nonlinear models [\S~\ref{parNonlinearLatticeOfParticles}].
Finally in \S~\ref{parAHypocoerciveSystemOfInteractingParticles} we will see
how one can consider time as a supplementary dimension of the system
to get contractivity results for non-reversible Markov chains
(\emph{hypocoercivity}) on an infinite system of particles.

\begin{NOTA}
In this chapter, all the probability systems considered
will be endowed with their natural $\sigma$-metalgebras,
cf.\ Definition~\ref{def3128}.
To alleviate notation, I will give no names to these $\sigma$-metalgebras,
but will plainly write ``${\{X:Y\}_*}$'' to mean
``the subjective decorrelation between~$X$ and~$Y$
seen from the natural $\sigma$-metalgebra of the underlying system''.
\end{NOTA}



\section{Back to Ising's model}\label{parBackToIsingsModel}

\subsection{Standard Ising's model}

In all this section, we work on the lattice $\ZZ^n$
equipped with its natural distance $\dist$;
accordingly $|\Bcdot|$ will denote the $\ell^1$ norm on~$\RR^n$.
Recall the definition of Ising's model and the related notation that we introduced
in \S~\ref{parSomeResultsOnIsingsModel}, and Theorem~\ref{t4505}
on the existence of a completely analytical regime.

The following theroem states that Ising's model in completely analytical regime
is $\rho$-mixing, i.e.\ that two distant bunches of spins
are little correlated in the sense of maximal correlation:
%
\begin{Thm}\label{t3730}
For Ising's model on~$\ZZ^n$ in the completely analytical regime,
\begin{ienumerate}
\item\label{i3743a} There exists some $\psi'>0$ (the same as in Theorem~\ref{t4505})
such that for all disjoint $I,J\subset\ZZ^n$, one has when~$\dist(I,J)\longto\infty$ that
\[\label{f4615}
\{ \vec\omega_I : \vec\omega_J \} \leq \exp\big[-\big(\psi'+o(1)\big) \, \dist(I,J)\big] ,\]
where the ``$o(1)$'' can be
easily computed as an explicit function of~$\dist(I,J)$, $n$, $T$, $\psi'$ and the
$C'$ appearing in Theorem~\ref{t4505}.
\item\label{i3743b} There exists some $k<1$ such that for all disjoint $I,J\subset\ZZ^n$,
\[ \{ \vec\omega_I : \vec\omega_J \} \leq k .\]
\item\label{i5202}
Points~(\ref{i3743a}) and~(\ref{i3743b}) remain valid uniformly under any law of the form
$\Pr[\Bcdot|{\vec\omega_K=\vechat{\omega}_K}]$,
for~$K \subset \ZZ^n$ and~$\vechat{\omega}_K \in \{\pm1\}^K$ a `boundary condition' on~$K$.
\end{ienumerate}
\end{Thm}

\begin{Rmk}
Let us compare Theorem~\ref{t3730} with Theorem~\ref{t4505}.
Both theorems state decorrelation between distant bunches of spins
above temperature $T_{\mrm{c}}'$; the difference relies in using $\rho$-mixing
rather than $\beta$-mixing to quantify dependence between the bunches
in Theorem~\ref{t3730}.

Both results give an exponential decay of correlations,
with the same exponential constant $\psi'$,
but Theorem~\ref{t3730} is more powerful
in the sense that the bound~(\ref{f4615}) is uniform in the size of~$I$ and~$J$
while (\ref{f7869}) was not.
Moreover, thanks to Point~(\ref{i3743b}) we get a non-trivial result
for any choice of disjoint~$I$ and~$J$, which was not the case beforehand.
Recall that the drawbacks of Theorem~\ref{t4505} were inherent to $\beta$-mixing,
as Theorem~\ref{c3793} shew.

Both result remain valid under conditioning.
However, if one takes a \emph{random} boundary condition%
—that is, if one works under the law $\Pr[\Bcdot|{\vec\omega_K\in C}]$
for some non-singleton $C \subset \{\pm1\}^K$\hbox{ —,}
then Point~(\ref{i5202}) of Theorem~\ref{t3730} fails (cf.\ Remark~\ref{r4670}),
while (\ref{f7869}) is still valid by convexity of the total variation norm.
\end{Rmk}

\begin{Rmk}
Let us compare Theorem~\ref{t3730} with Theorem~\ref{c6177}.
The result of Theorem~\ref{c6177} can be rewritten:
\[ \big\{ \vec\omega_{\{0\}\times\ZZ^{n-1}} : \vec\omega_{\{x\}\times\ZZ^{n-1}} \big\}
\leq e^{-\psi x} . \]
Theorem~\ref{t3730} can be seen as a generalization of that result
to the case where $I$ and~$J$ have arbitrary shapes.%
\footnote{Note that in the case $I$ and~$J$ are hyperplanes,
we shew on page~\pageref{parHyperplanesInIsingsModel}
that (\ref{f4615}) could be improved into
\[ \big\{ \vec\omega_{\{0\}\times\ZZ^{n-1}} : \vec\omega_{\{x\}\times\ZZ^{n-1}} \big\}
\leq e^{-\psi' x} .\]}
Moreover, Point~(\ref{i5202}) also gives the existence of a conditional version,
which we did not have before.

There is however a price to pay for this greater generality,
since we had to require complete analyticity rather than just weak mixing,
which can be really more restrictive in some cases
(cf.\ Footnote~\ref{Tc'=Tc?} on page~\pageref{Tc'=Tc?}).
\end{Rmk}

\begin{Rmk}
Continuing the previous remark,
a natural open question is whether one can tensorize maximal decorrelation
under assumptions of weak mixing type.
%
In the case of Ising's model at least, I expect $\rho$-mixing to remain true%
—even for arbitrary shapes—as soon as $T>T_{\mrm{c}}$,
because on the one hand Theorem~\ref{c6177}
proves $\rho$-mixing between parallel hyperplanes,
while on the other hand $\rho$-mixing
seems to hold also in the `opposite extreme case'
when $I$ and~$J$ make a check pattern.

By the way, it is likely that the natural condition
should not be weak mixing itself but rather something like \emph{strong mixing for cubes}
(often called merely \emph{strong mixing}%
\footnote{Strong mixing \emph{stricto sensu} is actually
the same as complete analyticity, so that mathematicians have got used to undermeaning
``for cubes''—but strong mixing for cubes is strictly weaker than complete analyticity!%
~\cite[\S~2]{MO}.},
which means that when a boundary condition is fixed outside a cube of arbitrary edge,
changing one spin on the boundary has an effect in total variation
which decreases exponentially with the distance to the spin changed.
%
In fact it has been proved~\cite{MOS} that in dimension~2,
weak mixing is equivalent to strong mixing.
\end{Rmk}

\begingroup\def\proofname{Proof of Theorem~\ref{t3730}}\begin{proof}
Theorem~\ref{t3730} will be a direct consequence of the work of Chapter~\ref{parTensorization}
as soon as we show that, denoting by~$*$ the natural $\sigma$-metalgebra of the system
(i.e.\ the $\sigma$-metalgebra generated by the~$\omega_i$), for all distinct $i,j\in\ZZ^n$,
one has
\[\label{f7295} \{ \omega_i : \omega_j \}_* \leq c_0C' e^{-\psi' \dist(i,j)} \wedge k_0 \]
for some explicit $c_0 < \infty$ and~$k_0 < 1$ only depending on~$n$ and~$T$.
Then indeed, Proposition~\ref{pro8882} yields
%
\begin{multline}
\{ \vec\omega_I : \vec\omega_J \} \leq
\sum_{\substack{\delta\in\ZZ^n\\|\delta|\geq\dist(I,J)}}
c_0C' e^{-\psi'|\delta|}
= c_0C' \sum_{d=\dist(I,J)}^\infty
\#\{\delta\in\ZZ^n\colon |\delta|=d\} e^{-\psi'd} \\
\stackrel{\mathsmaller{\dist(I,J)\longto\infty}}{\sim}
c_0C' \sum_{d=\dist(I,J)}^\infty \frac{2^nd^{n-1}}{(n-1)!} e^{-\psi'd}
\sim \frac{c_0C'2^n}{(n-1)!} \dist(I,J)^{n-1} e^{-\psi'\dist(I,J)} \\
= e^{-(\psi'+o(1))\dist(I,J)},
\end{multline}
whence Point~(\ref{i3743a}). Moreover, since
\[ \sum_{\substack{\delta\in\ZZ^n\setminus\{0\}
\\|\delta|\leq\dist(I,J)}} c_0C' e^{-\psi' \dist(i,j)}
< \infty ,\]
Point~(\ref{i3743b}) follows from Lemma~\ref{lem4067},
and finally (\ref{i5202}) is a consequence of \S~\ref{parSubjectiveResults}
about subjective results.

So, we have to prove~(\ref{f7295}).
Let~$\vechat{\omega}_K \in \{\pm1\}^K$, $K\subset\ZZ^n$, be some arbitrary boundary condition,
and denote by~$\Pr_{\mathrm{con}}$ the associated law,
that is, $\Pr_{\mathrm{con}} = \Pr[\Bcdot|{\vec\omega_K = \vechat{\omega}_K}]$;
our goal is to show that under~$\Pr_{\mathrm{con}}$, for all distinct $i,j\in\ZZ^n$, one has
$\{ \omega_i : \omega_j \} \leq c_0C' e^{-\psi' \dist(i,j)} \wedge k_0$.

The result is immediate if $i\in K$, resp.\ $j\in K$
(since then $\omega_i$, resp.~$\omega_j$, is constant and thus independent of everything),
so we assume $i,j\notin K$.
%
We begin with observing that if $K$ is the set~$N(i)$ of all the neighbours of~$i$,
equilibrium at~$i$ implies that, whatever the boundary condition may be:
\[\label{f7537} \Pr_{\mathrm{con}}[\omega_i = -1], \Pr_{\mathrm{con}}[\omega_i = +1]
\geq \big( e^{4n/T} + 1 \big)^{-1} \]
—the extremal cases being when $\vechat{\omega}_{N(i)} \equiv +1$,
resp.\ $\vechat{\omega}_{N(i)} \equiv -1$.
Now in the general case $K \subset \ZZ^n\setminus\{i\}$,
$\Law_{\mathrm{con}}[\omega_i]$ is an average of laws of the form
$\Law(\omega_i|{\vec\omega_{N(i)}=\vechat{\omega}_{N(i)}})$,
so that (\ref{f7537}) remains valid.
%
Similarly, equilibrium on~$\{i,j\}$ gives that
for all~$a,b\in\{\pm1\}$,
\[ \label{f7538} \Pr_{\mathrm{con}}[\omega_i = a \text{ and } \omega_j = b]
\geq \big( e^{8n/T} + 2e^{(4n+2)/T} + 1 \big)^{-1} .\]

Now, recall that the correlation level between two two-ranged variables
can be computed by Formula~(\ref{f0699}), where
$|p_a^b-p_ap^b|$ is also $\beta(X,Y)/2$.
Thus the bound ``$\{\omega_i:\omega_j\} \leq C_0 e^{-\psi' \dist(i,j)}$''
is a direct consequence of Theorem~\ref{t4505}, with
\[ c_0 = \frac{1/2}{(e^{4n/T}+1)^{-1}\big(1-(e^{4n/T} + 1)^{-1}\big)}
= \tanh(4n/T) + 1 .\]

It remains to prove the bound ``$\{\omega_i:\omega_j\} \leq k_0$''.
We will use the following corollary of~(\ref{f0699}):
%
\begin{Lem}\label{l1373}
With the notation of Remark~\ref{rmk9549},
there exists $a,b$ in the respective ranges of~$X,Y$ such that
\[ \{X:Y\} \leq 1 - 4p_a^b .\]
\end{Lem}

\begingroup\def\proofname{Proof of Lemma~\ref{l1373}}%
\begin{proof}
The difference~$p_a^b-p_ap^b$ gets its sign changed whenever $a$, resp.~$b$, changes,
so there are some~$a$ and~$b$ for which this value is nonpositive;
moreover, denoting by~$\{a,a'\}$ and~$\{b,b'\}$ the respective ranges of~$X$ and~$Y$,
$p_{a'}^{b'}-p_{a'}p^{b'}$ is also nonpositive.
Now one has
\[ \frac{p_ap^b}{\sqrt{p_ap_{a'}p^bp^{b'}}} \times \frac{p_{a'}p^{b'}}{\sqrt{p_ap_{a'}p^bp^{b'}}}
= 1 ,\]
so that either $p_ap^b$ or $p_{a'}p^{b'}$ is $\leq \sqrt{p_ap_{a'}p^bp^{b'}}$.
Up to changing notation we can assume that it is~$p_ap^b$,
and then
\[ \{X:Y\} = \frac{|p_a^b-p_ap^b|}{\sqrt{p_ap_{a'}p^bp^{b'}}}
= \frac{p_ap^b - p_a^b}{\sqrt{p_ap_{a'}p^bp^{b'}}}
\leq 1 - \frac{p_a^b}{\sqrt{p_ap_{a'}p^bp^{b'}}} \leq 1 - 4p_a^b .\]
\end{proof}\endgroup

Combining Lemma~\ref{l1373} with~(\ref{f7538}),
we then get the desired bound, with
\[ k_0 = 1 - 4\big( e^{8n/T} + 2e^{(4n+2)/T} + 1 \big)^{-1} < 1 .\]
\end{proof}\endgroup


Formula~(\ref{f7295}) is also what we need to apply the results of Chapter~%
\ref{parOtherApplicationsOfTensorizationTechniques}. Indeed, denoting
$\epsilon(z) \coloneqq \{ X_i : X_{i+z} \}_*$, it gives that $\sum_{z\in\ZZ^n} \epsilon(z) < \infty$
with $\epsilon(z) < 1$ as soon as $z \neq 0$,
so that Theorems~\ref{thm3015} and~\ref{thm7431} yield respectively:
%
\begin{Thm}
In completely analytical regime, the spins Ising's model satisfies
the central limit theorem, in the sense that the conclusions of Theorem~\ref{thm3015}
hold for them.
\end{Thm}
%
\begin{Thm}
In completely analytical regime, the Glauber dynamics for Ising's model
has a (strictly) positive spectral gap,
and this remains valid uniformly if one fixes
a `boundary condition' on the spins of some $K \subset \ZZ^n$.
\end{Thm}

\begin{Rmk}
As I told in Chapter~\ref{parOtherApplicationsOfTensorizationTechniques},
results of these kinds have already been studied by other methods
(see e.g.~\cite{Bradley92, Dedecker} for the CLT and~\cite{Martinelli} for the spectral gap).
For the standard Ising model in completely analytical regime,
which is ``very nice'', these previous works apply well,
so the two theorems above are not new.
They are interesting however because of the new method used to prove them,
which is quite direct and likely to apply to a broader class of models.
Such models will be presented in the sequel of this chapter.
\end{Rmk}


\subsection{Generalizations of Ising's model}

The previous results can be adapted to several kinds of generalizations of Ising's model.
Let us expose some of them.

\subsubsection*{Long-range Ising models}

A physically important case is the \emph{long-range Ising models} on~$\ZZ^n$.
In these models, the states space is unchanged, but the Hamiltonian $H$ becomes
\[ H(\vec\omega) = - \frac{1}{2} \sum_{i\neq j} J(j-i) \omega_i\omega_j ,\]
where $J \colon \ZZ^n\setminus\{0\} \to \RR$ is some symmetric function with non-compact support
such that $J(z) \stackrel{|z|\longto\infty}{=} O(|z|^{-(n+\alpha)})$ for some $\alpha>0$.

Let us state a decorrelation result for this class of models.
The frame of the proof of the following proposition will work as well
for the other generalizations of Ising's model.
%
\begin{Pro}\label{pro5088}
There exists an temperature $T_1 < \infty$ such that, provided $T\geq T_1$:
\begin{ienumerate}
\item\label{i0797} Equilibrium for the long-range Ising model is unique;
\item\label{i96a} Uniformly in~$i,j$,
$\{\omega_i:\omega_j\}_* \stackrel{|j-i|\longto\infty}{=} O(|j-i|^{-(n-\alpha)})$;
\item\label{i96b} There exists some $k_0 < 1$ such that for all~$i\neq j$,
$\{\omega_i:\omega_j\}_* \leq k_0$.
\end{ienumerate}
\end{Pro}

\begin{proof}
The principle of the proof consists in coupling two Glauber dynamics
with different initial conditions.
Recall that the Glauber dynamics is defined as follows:
each spin has an independent clock ringing with rate $1$,
and when the clock of a spin rings, this spin is flipped so that its final state is drawn
according to its equilibrium measure conditionnally to the state of all other spins.
Namely, if the clock of spin $i$ rings at time~$t$,
denoting as usual $\beta = T^{-1}$,
\[ \Pr[\omega_i(t+)=+1] = \frac{\exp\big(\beta\sum_{j\neq i}J(j-i)\omega_j(t)\big)}
{2\cosh\big(\beta\sum_{j\neq i}J(j-i)\omega_j(t)\big)} \]
and $\Pr[\omega_i(t+)=-1] = 1-\Pr[\omega_i(t+)=+1]$.

To couple the Glauber dynamics, we will assume that, rather than just ``ringing'' at time~$t$,
the clock of~$i$ is a Poisson process on~$\RR_+ \times (0,1)$,
points of which are denoted by~$(t,y)$.
Then, if at time~$t$ the clock of spin $i$ has a point $(t,y)$,
spin $i$ flips to~$+1$ if $y<\Pr[\omega_i(t+)=+1]$,
resp.\ to~$-1$ if~$y\geq\Pr[\omega_i(t+)=+1]$.

Now, consider two Glauber dynamics~$\vec\omega^-$ and~$\vec\omega^+$
having the same Poisson process, but starting with different initial conditions.
It will be convenient%
\footnote{In the cases where interactions can be antiferromagnetic ($J<0$),
monotonicity does not stand any more;
the proof however remains valid with a heavier formalism,
replacing ``$>$'' by ``$\neq$'' and putting absolute values at the right places.}
to assume that $\vec\omega^-(t=0) \leq \vec\omega^+(t=0)$ almost-surely:
then, as we will see, for the coupled dynamics
one has (a.s.) $\vec\omega^-(t) \leq \vec\omega^+(t) \ \forall t$.
At time $t$, denote by~$\Theta(t)$ the set of points where~$\vec\omega^-$ and~$\vec\omega^+$ differ:
\[ \Theta(t) = \big\{ i\in\ZZ^n \colon \big(\omega^-_i(t), \omega^+_i(t)\big) = (-1,+1) \} .\]
%
When the clock at spin $i$ rings at time $t$, three cases have to be distinguished:
\begin{enumerate}
\item If $y < \exp\big(\beta\sum_{j\neq i}J(j-i)\omega^-_j(t)\big) \mathbin{\big/}
2\cosh\big(\beta\sum_{j\neq i}J(j-i)\omega^-_j(t)\big)$,
then both $\omega^+_i$ and~$\omega^-_i$ flip into state~$+1$;
\item If $y > \exp\big(\beta\sum_{j\neq i}J(j-i)\omega^+_j(t)\big) \mathbin{\big/}
2\cosh\big(\beta\sum_{j\neq i}J(j-i)\omega^+_j(t)\big)$,
then both $\omega^+_i$ and~$\omega^-_i$ flip into state~$-1$;
\item If $\frac{\exp\big(\beta\sum_{j\neq i}J(j-i)\omega^-_j(t)\big)}
{2\cosh\big(\beta\sum_{j\neq i}J(j-i)\omega^-_j(t)\big)} < y
< \frac{\exp\big(\beta\sum_{j\neq i}J(j-i)\omega^+_j(t)\big)}
{2\cosh\big(\beta\sum_{j\neq i}J(j-i)\omega^+_j(t)\big)}$,
then $\omega^+_i$ flips into state $+1$ while $\omega^-_i$ flips into state $-1$.
\end{enumerate}
%
Denoting
\[ \mcal{J} \coloneqq \sum_{z\in\ZZ^n\setminus\{0\}} J(z) ,\]
which is always finite by the assumption on~$J$,
the probability of each of the two first cases is bounded below by
$e^{-\beta\mcal{J}} / 2\cosh(\beta\mcal{J})$.
%
The probability of the third case is
\[ \frac{\sinh\big(2\beta\sum_{j\in\Theta(t)\setminus\{i\}} J(j-i)\big)}
{2\cosh\big(\beta\sum_{j\neq i}J(j-i)\omega^-_j(t)\big)
\cosh\big(\beta\sum_{j\neq i}J(j-i)\omega^+_j(t)\big)} ,\]
which is bounded above by~$\beta \sum_{j\in\Theta(t)\setminus\{i\}} {J(j-i)}$
thanks to the following computational
%
\begin{Lem}
For~$a \leq b$ two real numbers,
\[ \sinh(b-a) \leq (b-a) \cosh a \cosh b .\]
\end{Lem}
%
\begin{proof}
Making the change of variables $x = (a+b)/2, t = (b-a)/2$,
we have to prove that for~${x \in \RR}, \ab t \geq 0$, one has:
\[\label{fo4994} \sinh(2t) \leq 2 t \cosh(x-t)\cosh(x+t) .\]
%
If we consider the right-hand side of~(\ref{fo4994})
as a function of~$x$, it is symmetric (since $\cosh$ is symmetric)
and its logarithm is convex (since $\ln\circ\cosh$ is convex,
its derivative being the increasing function~$\tanh$),
so its minimum is attained for $x=0$;
thus it suffices to prove~(\ref{fo4994}) in that case,
i.e.\ to prove that $\sinh(2t) \leq 2t\cosh^2t$ for all~$t\geq 0$.
But $\sinh(2t) = 2\sinh t\cosh t$, so we can simplify both sides by~$2\cosh t$,
and then it suffices to prove that
$\sinh t \leq t \cosh t$, which is true since
$\tanh t \leq t$ for all~$t\geq 0$.
\end{proof}

Thanks to these estimates, we can define a process Markovian $\Theta^*(t)$ on~$\mathfrak{P}(\ZZ^n)$
such that almost-surely, $\Theta^*(t) \supset \Theta(t)\ \forall t$.
This process has the following law:
%
\begin{Def}
The law of~$\Theta^*$ is defined thanks to independent Poissonian clocks
indexed by~$(\ZZ^n)^2$.
For $i\neq j$ the clock $(i,j)$ has rate $\beta J(j-i)$,
while the clock $(i,i)$ has rate $e^{-\beta\mcal{J}} / \cosh(\beta\mcal{J})$.
At $t=0$ one has $\Theta^*(0) = \Theta(0)$.
If at time $t$ the clock $(i,j)$ rings, with $j\neq i$, then:
%
\begin{itemize}
\item Either $i\in\Theta^*(t-)$ and then $\Theta^*$ changes so that
$\Theta^*(t+) = \Theta^*(t-) \cup \{j\}$%
\footnote{Of course, if $j\in\Theta^*(t-)$ then $\Theta^*$ does actually not change.};
\item Or $i\notin\Theta^*(t-)$ and then $\Theta^*$ does not change.
\end{itemize}
On the other hand, if at time $t$ the clock $(i,i)$ rings,
then $\Theta^*$ changes so that $\Theta^*(t+) = \Theta^*(t-) \setminus \{i\}$.
\end{Def}

Let~$\lambda \coloneqq \beta\mcal{J} - \big(e^{-\beta\mcal{J}}/\cosh(\beta\mcal{J})\big)$.
If we take $\EE\big[\#\Theta(t=0)\big] < \infty$%
\footnote{The general case where $\Theta^*$ can be infinite can be got
from the finite case by passing to the limit,
despite some technicalities of little interest.},
it is immediate that $\#\Theta^*(t)/e^{\lambda t}$ is a supermartingale.
So, provided $T$ is large enough so that $\lambda < 0$, i.e.
\[\label{for5229} \beta\mcal{J} < \frac{e^{-\beta\mcal{J}}}{\cosh(\beta\mcal{J})} ,\]
the two processes~$\vec{\omega}^-(t)$ and~$\vec{\omega}^+(t)$
tend to be equal when~$t\longto\infty$; in particular they have the same equilibrium.
That proves Point~(\ref{i0797}) of the Lemma,
since any initial condition stands between the `extreme' conditions
$\vec\omega^-(t=0) \equiv -1$ and $\vec\omega^+(t=0) \equiv +1$.

Observe that the previous reasoning remains entirely valid if one reasons conditionally
to some boundary condition of the form ``$\vec\omega_K = \vechat{\omega}_K$'',
with the same condition on~$T$.

Now we are turning to the correlation between two distant spins.
Let~$i\in\ZZ^n$ and let~$\vechat{\omega}_K$ be some boundary condition on
some $K \subset \ZZ^n\setminus\{i\}$.
Suppose $T$ satisfies~(\ref{for5229});
I want to compare the Glauber dynamics corresponding to the boundary condition
``$\vec\omega_{K\uplus\{i\}} = \big(\vechat{\omega}_K, (+1)^{\{i\}}\big)$''%
—where $(\vechat{\omega}_K, (+1)^{\{i\}})$ stands for the function on
$K\uplus\{i\}$ which is equal to~$\hat\omega$ on~$K$ and to~$+1$ at~$i$ —%
with the Glauber dynamics corresponding to the boundary condition
``$\vec\omega_{K\uplus\{i\}} = \big(\vechat{\omega}_K, (-1)^{\{i\}}\big)$''.
In this frame, one defines the process $\Theta^*$ as previously,
except that one imposes that $\Theta^*(t) \cap K = 0$ and $i \in \Theta^*(t)$ for all~$t$.
This time, it is the equilibrium behaviour of~$\Theta^*$ which interests us.
Denote by~$\Pr_{\text{eq}}$ the equilibrium law of~$\Theta^*$;
for~$j'\in\ZZ^n\setminus K$, denote $\theta(j) \coloneqq \Pr_{\text{eq}}[j\in \Theta^*]$.
Then $\theta$ satisfies the following discrete subelliptic equation
with Dirichlet boundary conditions:
%
\[\label{fo7960} \left\{ \begin{array}{ll}
\forall j\notin K\uplus\{i\}
\quad \frac{e^{-\beta\mcal{J}}}{\cosh(\beta\mcal{J})} \theta(j) \leq
\beta \sum_{\substack{i'\in\ZZ^n\setminus K \\ i'\neq j}} J(i'-j) \theta(i') ;\\
\forall k\in K \quad \theta(k) = 0; \qquad \theta(i) = 1.
\end{array} \right. \]

Define the convolution kernel $a$ on~$\ZZ^n$ by
\[ \left\{ \begin{array}{ll} a(0) = 1 ; \\
\forall z\neq 0 \quad a(z) = - \frac{\cosh(\beta\mcal{J})}{e^{-\beta\mcal{J}}}\beta J(z) ,
\end{array}\right. \]
so that (\ref{fo7960}) writes in the bulk:
\[ a * \theta \leq 0 .\]
%
Writing $a \eqqcolon \delta_0 - \tilde{a}$,
Condition~(\ref{for5229}) ensures that $\|\tilde{a}\|_{\ell^1} < 1$.
Since $\ell^1(\ZZ^n)$ is a Banach algebra for the convolution operator $*$,
with neutral element $\delta_{0}$, it follows that $a$ is invertible with inverse
\[\label{fo7178} a^{-*} =
\delta_{0} + \tilde{a} + \tilde{a} * \tilde{a} + \tilde{a} * \tilde{a} * \tilde{a} + \cdots .\]
Since $\tilde{a} \geq 0$, $a^{-*}$ is nonnegative everywhere with $a^{-*}(0)>0$.
Therefore the function~$F \coloneqq \big(a^{-*}(0)\big)^{-1} \delta_i * a^{-*}$
satisfies:
\[\label{fo7960*} \left\{ \begin{array}{ll}
\forall j\notin K\uplus\{i\} \quad \big(a \ast F\big)(j) = 0 ;\\
\forall k\in K \quad F(k) \geq 0; \qquad F(i) = 1.
\end{array}\right.\]
Comparing~(\ref{fo7960}) with~(\ref{fo7960*}),
since~(\ref{fo7960}) is subelliptic,
we can apply a maximum principle to it%
\footnote{The maximum principle is generally stated in a PDE context,
see for instance~\cite[\S~3.1]{elliptique}, but it works exactly the same for discrete equations.},
which yields that $\theta \leq F$ everywhere.
But $J(z) = O(|z|^{-(n+\alpha)})$,
so by Lemma~\ref{l1939b} in appendix, $F(j) = O(|j-i|^{-(n+\alpha)})$,
and therefore
\[\label{fo9696} \Pr\big[ \omega_j=+1 \big| \vec\omega_K = \vechat{\omega}_K , \omega_i = 1 \big]
- \Pr\big[ \omega_j=-1 \big| \vec\omega_K = \vechat{\omega}_K , \omega_i = -1 \big]
= O(|j-i|^{-(n+\alpha)}) ,\]
uniformly in~$i, j, K, \vechat{\omega}_K$.

The end of the proof, namely deducing Point~(\ref{i96a}) from~(\ref{fo9696})
and proving Point~(\ref{i96b}), is then performed in the same way as to establish~(\ref{f7295})
in the proof of Theorem~\ref{t3730}.
\end{proof}

Thanks to Proposition~\ref{pro5088}, we can apply the results of
Chapters~\ref{parTensorization} and~\ref{parOtherApplicationsOfTensorizationTechniques}.
One gets the following
%
\begin{Thm}\label{thm8777}
For the long-range Ising model on~$\ZZ^n$ at~$T\geq T_1$,
\begin{ienumerate}
\item\label{i204a} For all disjoint $I,J\subset\ZZ^n$, uniformly in~$I,J$, one has an estimate
\[\label{fo2580}
\{ \vec\omega_I : \vec\omega_J \} \leq O \big(\dist(I,J)^{-\alpha}\big) ,\]
where the~$O(\Bcdot)$ can be turned into an explicit constant only depending
on~$\mcal{J}$ and~$T$.
Moreover, there exists some $k<1$
(still explicit and only depending on~$\mcal{J}$ and~$T$)
such that for all disjoint $I,J\subset\ZZ^n$,
\[ \{ \vec\omega_I : \vec\omega_J \} \leq k .\]
%
\item\label{i204b} The spins satisfies
the central limit theorem, in the sense that the conclusions of Theorem~\ref{thm3015}
hold for them.
%
\item\label{i204c} The Glauber dynamics has a positive spectral gap.
%
\item\label{i204d}
Points~(\ref{i204a}) and~(\ref{i204c}) remain valid uniformly under any law of the form
$\Pr[\Bcdot|{\vec\omega_K=\vechat{\omega}_K}]$,
for~$K \subset \ZZ^n$ and~$\vechat{\omega}_K \in \{\pm1\}^K$ a `boundary condition' on~$K$.
\end{ienumerate}
\end{Thm}
%
\begin{proof}
The proof is the same as the work done in the previous subsection.
The only difference is to prove~(\ref{fo2580}),
which follows from the following computation:
denoting $D\coloneqq\dist(I,J)$, one has that, when~$D\longto\infty$,
\[ \sum_{\substack{z\in\ZZ^n\\|z|\geq D}} \frac{1}{|z|^{n+\alpha}}
\leq \sum_{d=D}^\infty
\frac{\#\{z\in\ZZ^n\colon |z|=d\}}{d^{n+\alpha}}
= \sum_{d=D}^\infty \frac{O(d^{n-1})}{|d^{n+\alpha}|}
= O\Big( \sum_{d=D}^\infty \frac{1}{d^{1+\alpha}} \Big) = O(D^{-\alpha}). \]
\end{proof}


\subsubsection*{Spin glasses}

\emph{Spin glasses} are another generalization of Ising's model.
In these models, the interaction constants are not invariant by translation any longer.
The Hamiltonian writes
\[ H(\vec\omega) = - \frac{1}{2} \sum_{i\neq j} J(i,j) \omega_i\omega_j \]
(with $J(j,i) = J(i,j)$), where the~$J(i,j)$ themselves are random.
We make the following assumptions on the interaction constants:
%
\begin{Hyp}\label{hyp79}
For distinct unordered pairs $\{i,j\}$, all the~$J(i,j)$ are independent.
Moreover, $J(i,j)$ is distributed according to some law $P_J^{(j-i)}$ only depending on~$(j-i)$%
\footnote{Observe that one has necessarily $P_J^{(-z)} = P_J^{(z)}$ for all~$z$;
in particular the function~${J_\infty \colon} \ab {\ZZ^n \setminus \{0\}} \to \RR_+$
shall always be symmetric.}.
We will assume that all the~$P_J^{(z)}$ have bounded support,
and we denote by~$J_\infty(z)$ the smallest number such that
$P_J^{(z)} \big[ |J| \leq J_\infty(z) \big] = 1$.
\end{Hyp}

\begin{Rmk}
Here the~$J(i,j)$ can be negative, which corresponds to antiferromagnetic interactions.
\end{Rmk}

\begin{NOTA}
In spin glass models, there are two levels of randomness: first to fix the~$J(i,j)$,
next to take $\vec\omega$ according to the Gibbs measure associated to~$H$.
When both levels of randomness are taken into consideration, one speaks of \emph{annealed} law.
Here I am only interested in the \emph{quenched} laws,
which deal with the second level of randomness for fixed~$J(i,j)$.
I will write sentences beginning with ``for almost-all quenched systems'',
which mean that what follows is valid for almost-all Gibbs measures
when the~$J(i,j)$ are taken randomly according to Assumption~\ref{hyp79}.
\end{NOTA}

The machinery exposed above still works for spin glass models.
We obtain the
%
\begin{Thm}
Suppose that when~$|z| \longto \infty$,
$J_\infty(z)$ decreases at least as fast as $O(|z|^{-(n+\alpha)})$ for some $\alpha>0$.
Then there is a $T_1 < \infty$ such that, for the spin glass model on~$\ZZ^n$ at~$T\geq T_1$,
for almost-all quenched systems,
\begin{ienumerate}
\item If $J_\infty(z) = O(|z|^{-(n+\alpha)})$,
then for all disjoint $I,J\subset\ZZ^n$, uniformly in~$I,J$, one has an estimate
\[\label{fo8727} \{ \vec\omega_I : \vec\omega_J \} \leq O \big(\dist(I,J)^{-\alpha}\big) .\]
%
If moreover $J_\infty(z)$ has exponential decay
(see Definition~\ref{def726} in the appendix),
then the right-hand side of~(\ref{fo8727}) can even be replaced by
``$\theta\big(\dist(I,J)\big)$'' for some function~$\theta(\Bcdot)$ with exponential decay.

In both cases, there exists some $k<1$ such that for all disjoint $I,J\subset\ZZ^n$,
\[ \{ \vec\omega_I : \vec\omega_J \} \leq k .\]
%
\item Points~(\ref{i204b})--(\ref{i204d}) of Theorem~\ref{thm8777} hold.
\end{ienumerate}
\end{Thm}


\subsection*{Synthetic vocabulary}

For all the models considered in this section,
the techniques used and the results stated walked along the same lines.
First, one establishes a bound $\{X_i:X_j\}_* \leq \epsilon(j-i) \wedge k_0$ for all~$i\neq j$,
for some sufficiently rapidly decreasing function
$\epsilon \colon \ZZ^n \to [0,1]$ and some $k_0 < 1$.
Then, one applies the results of
Chapters~\ref{parTensorization} and~\ref{parOtherApplicationsOfTensorizationTechniques},
which yield maximal decorrelation for distant bunches of spins
(which is sometimes called \emph{(interlaced) $\rho^*$-mixing})
with uniformly non-full correlation between any two disjoint bunches of spins
(which is sometimes denoted ``$\rho^*(1)<1$''),
central limit theorem, and spectral gap for the Glauber dynamics.

Since this method will be used again in the following sections,
it will be convenient to introduce some synthetic vocabulary:
%
\begin{Def}\label{def-WRM}
If a spin model (spins can have arbitrary range) $\vec{X}$ on~$\ZZ^n$ satisfies some bound
``${\{X_i:X_j\}_*} \leq \epsilon(j-i) \wedge k_0$'' for all distinct $i,j \in \ZZ^n$,
with $\sum_{z\in\ZZ^n\setminus\{0\}} \epsilon(z) < \infty$ and $k_0 < 1$,
we say that this model is \emph{well-$\rho$-mixing}.
According to our results, for such a model one has $\rho^*$-mixing with $\rho^*(1)<1$,
CLT and spectral gap.

Moreover,
\begin{ienumerate}
\item If $\epsilon(z) = O(|z|^{-(n+\alpha)})$ when~$|z|\longto\infty$,
then we say that the model is \emph{$\alpha$-polynomially $\rho$-mixing}.
According to our results, in this case $\rho^*$-mixing is polynomial with rate $\alpha$,
i.e.\ Formula~(\ref{fo2580}) holds.
\item If $\epsilon(z)$ has exponential decay (cf.\ Definition~\ref{def726}),
then we say that the model is \emph{exponentially $\rho$-mixing}.
According to our results, in this case $\rho^*$-mixing has an exponential speed of decay
(but not with the same rate as $\epsilon(\Bcdot)$, cf.\ Remark~\ref{rmk9890}),
i.e.\ a formula similar to~(\ref{f4615}) holds.
\end{ienumerate}
\end{Def}


\section{Quadratic models}\label{parQuadraticModels}

\begin{NOTA}
In this subsubsection, an arbitrary norm $|\Bcdot|$ on~$\ZZ^n$ is fixed.
\end{NOTA}

\begin{Def}
In our \emph{quadratic model}, the states space is $\Omega = \RR^{\ZZ^n}$
for some $n\in\NN^*$. For~$\vec\omega_{\ZZ^n} \in \Omega$, $i\in\ZZ^n$,
the real number $\omega_i$ will be called the \emph{polarization of particle~$i$}.
Each particle $i$ is submitted to two types of forces:
%
\begin{itemize}
\item A \emph{pinning force}, preventing the particle from having a too large polarization,
which derives from the quadratic potential $\omega_i^2/2$;
\item \emph{Interaction forces}: each particle $j\neq i$ exerts a force on~$i$
which tends to make the polarizations of particles~$i$ and~$j$ equal;
this force derives from a quadratic potential $\gamma_{j-i} (\omega_j-\omega_i)^2/2$.
\end{itemize}
%
In other words, the Hamiltonian of the system is formally defined by
\[ H(\vec\omega) = \frac{1}{2}\sum_{i\in\ZZ^n} \omega_i^2
+ \frac{1}{4}\sum_{i\neq j} \gamma_{j-i} (\omega_j-\omega_i)^2 ,\]
where the~$\gamma_z$, for~$z\in\ZZ^n\setminus\{0\}$,
are nonnegative numbers which we impose to satisfy the symmetry condition
$\gamma_z = \gamma_{-z}$ for all~$z$.
%
Moreover we impose the that the sum of the~$\gamma_z$ is convergent, and we denote
\[\label{f2380} \Gamma \coloneqq \sum_{z\in\ZZ^n\setminus\{0\}} \gamma_z < \infty .\]
\end{Def}

The Hamiltonian $H$ is a quadratic function of~$\vec\omega$,
so at fixed parameter $\beta$ the (infinite-dimensional) random vector~$\vec\omega$
will be Gaussian (and centered).
Let us compute its covariance:
the probability density of~$\vec\omega$ w.r.t.\ the `Lebesgue measure' on~$\Omega$
is formally defined by
\[ \frac{\dx\Pr_\beta [\vec\omega]}{\prod_{i\in\ZZ^n}\dx{\omega_i}} \propto
\exp \Big({\textstyle\frac{1}{2}} \T{\omega} (\beta Q) \omega \Big), \]
where $Q$ is the (infinite-dimensional) symmetric matrix defined by
\[\label{f1243} \left\{ \begin{array}{ccl} Q_{ij} \coloneqq -\gamma_{j-i} & \quad & \text{for $i\neq j$;} \\
Q_{ii} \coloneqq 1 + \Gamma & \quad & \text{on the diagonal,} \end{array} \right.\]
%
thus the covariance matrix of~$\vec\omega$ is $(\beta Q)^{-1}$.
So we have to compute $Q^{-1}$, the inverse matrix of~$Q$.
Since $Q$ is a Toeplitz matrix (with $n$-dimensional indexes)%
\footnote{Recall that saying that matrix $Q$ is \emph{Toeplitz}
means that its entries $Q_{ij}$ only depend on~$(j-i)$.},
$Q^{-1}$ —if it exists—will be of the same form.
Now, knowing that it is a Toeplitz matrix, $Q$ is described by the function
$a_Q \colon \ZZ^n \to \RR$ such that for all~$i,j$, $Q_{ij} = a_Q(j-i)$.
With this notation, (\ref{f1243}) rewrites:
\[ \forall z\in\ZZ^n \qquad a_Q(z) = \1{z=0} (1+\Gamma) - \1{z\neq0} \gamma_z .\]
%
When coded by functions like $a_Q$,
the multiplication of Toeplitz matrices becomes the convolution product:
\[ \forall M, N \ \text{Toeplitz} \qquad a_{MN} = a_M * a_N .\]
So, $Q^{-1}$ will be the Toeplitz matrix
whose $a_{Q^{-1}}$ is the inverse of~$a_Q$ for the convolution product.
Thanks to Condition~(\ref{f2380}), such an inverse always exists:
indeed we can write $a_Q = (1+\Gamma) (\delta_{0} - \tilde{a}_Q)$,
where $\tilde{a}_Q$ is a nonnegative function with
$\|\tilde{a}_Q\|_{\ell^1} = \Gamma \div (1+\Gamma) < 1$,
so that $a_Q$ is invertible with
\[\label{f0024} a_Q^{-*} = (1+\Gamma)^{-1} \big(\delta_{0} + \tilde{a}_Q +
\tilde{a}_Q * \tilde{a}_Q + \tilde{a}_Q * \tilde{a}_Q * \tilde{a}_Q + \cdots \big) .\]
In the end, at parameter $\beta>0$ the covariance matrix of~$\vec\omega$ has entries:
\[\label{f3366} \Cov(\omega_i,\omega_j) = \frac{a_{Q^{-1}}(j-i)}{\beta} .\]

\begin{Rmk}
All the entries of~$\Cov(\vec\omega)$ are nonnegative,
which reflects the fact that all the interaction forces are attractive.
\end{Rmk}

\begin{Rmk}
Since $\Cov(\vec\omega)$ depends on~$\beta$ only through a constant factor,
the behaviour of the system is exactly the same, up to a multiplicative constant,
for all~$\beta > 0$. Hence the study of correlations will not depend on~$\beta$.
\end{Rmk}

\begin{NOTA}
In the sequel, we fix arbitrarily $\beta = 1$ and we denote~$\Pr$ for~$\Pr_{\beta=1}$.
\end{NOTA}

Since the model is Gaussian, by~(\ref{f3366}) and Theorem~\ref{pro1857}
one has for all~$i \neq j$:
\[ \{ \omega_i : \omega_j \} = \frac{a_{Q^{-1}}(j-i)}{a_{Q^{-1}}(0)} .\]
Now we have the following claim, with an immediate key corollary:
\begin{Clm}\label{l3475}
For all~$i\neq j$, for all~$K \subset \ZZ\setminus\{i,j\}$,
\[ \{ \omega_i : \omega_j \}_{\vec\omega_K} \leq \{ \omega_i : \omega_j \} .\]
\end{Clm}

\begin{Cor}\label{c6834}
Denoting by~$*$ the natural $\sigma$-metalgebra of the system, for all~$i\neq j$,
\[\label{fc6834} \{ \omega_i : \omega_j \}_* = \{ \omega_i : \omega_j \}
= \frac{a_{Q^{-1}}(j-i)}{a_{Q^{-1}}(0)} .\]
\end{Cor}

\begin{proof}
The proof of Claim~\ref{l3475} relies on the following claims:
\begin{Clm}\label{c4577a}
Up to an additive constant, $\Law(\vec\omega_{\ZZ^n}|\vec\omega_K=\vechat{\omega}_K)$
is the same for all~$\vechat{\omega}_K \in \RR^K$, i.e.\ there exists a vector-valued function
$\vechat{\omega}_K \mapsto \offset(\vechat{\omega}_K) \in \RR^{\ZZ^n}$
such that the law of~$\vec\omega_{\ZZ^n}$ under~$\Pr[\Bcdot|\ab {\vec\omega_K=\vechat{\omega}_K}]$
is the same as the law of~$\vec\omega_{\ZZ^n} + \offset(\vechat{\omega}_K)$
under~$\Pr[\Bcdot|\ab {\vec\omega_K \equiv 0}]$.
\end{Clm}

\begin{Lem}\label{c4577b}
For~$(X,Y)$ a two-dimensional centered Gaussian vector with~$X$ and~$Y$ non-degenerate,
\[\label{f5223} \{X:Y\} = \sqrt{\frac{\EE[X^2]}{\EE[Y^2]}} \, \big|\EE[Y|X=1]\big| .\]
\end{Lem}

\begin{Clm}\label{c4577c}
For~$K\subset\ZZ^n$, the function~$\offset$
defined in Claim~\ref{c4577a} is nondecreasing,
in the sense that each of the entries of~$\offset(\vechat{\omega}_K)$
is a nondecreasing function of each~$\hat\omega_k$ for~$k\in K$.
\end{Clm}

\begin{Clm}\label{c4577d}
For~$i\in\ZZ^n$, $K\subset\ZZ^n\setminus\{i\}$:
\[ \offset\big(1^{\{i\}},0^K\big) \leq \offset\big(1^{\{i\}}\big) ,\]
where $(1^{\{i\}},0^K)$ stands for the function on
$K\uplus\{i\}$ which is equal to~$1$ at~$i$ and to~$0$ on~$K$,
resp.~$1^{\{i\}}$ stands for the function on~$\{i\}$ mapping $i$ to~$1$.
\end{Clm}

Admit temporarily the claims.
Let~$i,j$ be distinct points of~$\ZZ^n$, let~$K\subset\ZZ^n\setminus\{i,j\}$
and let~$\vechat{\omega}_K \in \RR^K$; our goal is to compute
$\{ \omega_i : \omega_j \}$ under~$\Pr[\Bcdot|\ab \vec\omega_K=\vechat{\omega}_K]$.
First, by Claim~\ref{c4577a} we can suppose that $\vechat{\omega}_K \equiv 0$.
Now under~$\Pr[\Bcdot|\ab \vec\omega_K\equiv0]$, $(\omega_i,\omega_j)$ is still Gaussian
by the properties of Gaussian vectors, and it is centered by symmetry,
therefore by Lemma~\ref{c4577b}, $\{\omega_i:\omega_j\}$ is equal to
\[\label{f5276}
\sqrt{\frac{\EE[\omega_i^2|\vec\omega_K\equiv0]}{\EE[\omega_j^2|\vec\omega_K\equiv0]}} \, \big| \EE\big[\omega_j\big|\vec\omega_K\equiv0\enskip\text{and}\enskip \omega_i=1\big] \big|
= \sqrt{\frac{\EE[\omega_i^2|\vec\omega_K\equiv0]}{\EE[\omega_j^2|\vec\omega_K\equiv0]}} \, \big(\offset(1^{\{i\}},0^K)\cdot j\big) \]
—one has indeed $\offset(1^{\{i\}},0^K)\cdot j \geq 0$,
since by Claim~\ref{c4577c}, $\offset(1^{\{i\}},0^K) \geq \offset(0^{\{i\}\uplus K}) \equiv 0$.

Now, taking $K=\emptyset$ in~(\ref{f5276}), we find that under the law $\Pr$:
\[ \{\omega_i:\omega_j\} = \sqrt{\frac{\EE[\omega_i^2]}{\EE[\omega_j^2]}} \,
\big(\offset(1^{\{i\}})\cdot j\big) ,\]
which is $\geq \sqrt{\EE[\omega_i^2] \div \EE[\omega_j^2]} \, \big(\offset(1^{\{i\}},0^K)\cdot j\big)$
by Claim~\ref{c4577d}.
But up to switching the roles of~$i$ and~$j$, we can assume that
$\EE[\omega_i^2] \div \EE[\omega_j^2] \geq
\EE[\omega_i^2|\vec\omega_K\equiv0] \div \EE[\omega_j^2|\vec\omega_K\equiv0]$,
thus getting the desired result:
\[ \{\omega_i:\omega_j\}
\geq \sqrt{\frac{\EE[\omega_i^2|\vec\omega_K\equiv0]}{\EE[\omega_j^2|\vec\omega_K\equiv0]}} \, \big(\offset(1^{\{i\}},0^K)\cdot j\big) = \{\omega_i:\omega_j\}_{\vec\omega_K} .\]
\end{proof}

\begingroup\def\proofname{Proof of the claims}
\begin{proof}\strut\\
\indent\textit{Claim~\ref{c4577a}\ --}
It is a well-known property of Gaussian vectors,
which here is stated in an infinite-dimensional setting.

\indent\textit{Claim~\ref{c4577b}\ --}
Since $(X,Y)$ is centered Gaussian,
$Y^{\sigma(X)}$ is the orthogonal projection of the $L^2$ variable $Y$ on~$\RR X$,
so $\EE[Y|X=x] \propto x$.
Thus one has:
\[\label{c5212} \EE[XY] = \int x\, \EE[Y|X=x]\, \dx\Pr[X=x] \\
= \int x^2\, \EE[Y|X=1]\, \dx\Pr[X=x] = \EE[Y|X=1] \, \EE[X^2] .\]
But for such a Gaussian vector, Theorem~\ref{pro1857} gives that
\[ \{X:Y\} = \frac{|\Cov(X,Y)|}{\ecty(x)\ecty(y)} = \frac{|\EE[XY]|}{\sqrt{\EE[X^2]\EE[Y^2]}} ,\]
which combined with~(\ref{c5212}) gives~(\ref{f5223}).

\indent\textit{Claim~\ref{c4577c}\ --} First, notice that
$\EE[\vec\omega_{\ZZ^n}|\vec\omega_K\equiv0] = \vec0$, so that
\[\offset(\vechat{\omega}_K) =
\EE\big[\vec\omega \big| \vec\omega_K=\vechat{\omega}_K \big] .\]
Now, allowing temporarily $\beta$ to vary again, by the properties of Gaussian vectors,
the vector-valued variable $\EE_\beta[\vec\omega_{\ZZ^n}|\vec\omega_K=\vechat{\omega}_K]$
is Gaussian with constant expectation and covariance matrix proportional to~$\beta$.
Therefore, the common expectation of all these laws
is equal to the constant value of~$\vec\omega_{\ZZ^n}$ for $\beta=0$,
which is the~$\vec\omega$ minimising $H$ under the constraint ``$\vec\omega_K=\vechat{\omega}_K$'':
\[ \offset(\vechat{\omega}_K) =
\argmin_{\vec\omega_K=\vechat{\omega}_K} H(\vec\omega) .\]
%
Since it minimizes energy, the state $\offset(\vec\omega_K)$ is at equilibrium outside $K$.
In other words, it is the solution of the following subelliptic system:
\[\label{f4463} \left\{ \begin{array}{ll}
\forall i\in\ZZ^n\setminus K
\quad -\omega_i + \sum_{j\neq i} \gamma_{j-i}(\omega_j-\omega_i) = 0 ;\\
\forall i\in K \quad \omega_i = \hat\omega_i .\end{array} \right.\]
(That system is clearly subelliptic
because the pinning and interaction forces all are attractive).
By the maximum principle,
the solution of~(\ref{f4463}) is an increasing function of the boundary condition,
which was our claim.

\indent\textit{Claim~\ref{c4577d}\ --} Denote
$\vec\omega^{1}_{\ZZ^n} = \offset(1^{\{i\}})$.
Since $\offset(0^{\{i\}}) = 0^{\ZZ^n}$, by Claim~\ref{c4577c}
one has ${\vec\omega^1_{\ZZ^n} \geq 0^{\ZZ^n}}$.
Since obviously $\vec\omega^1_i = 1$, one has even
$\vec\omega^1_{\ZZ^n} \geq (1^{\{i\}},0^{\ZZ^n\setminus\{i\}})$.
In particular, $\vec\omega^1_{\{i\}\uplus K} \geq (1^{\{i\}},0^K)$;
therefore, using again Claim~\ref{c4577c},
\[\label{f3643} \offset\big(\vec\omega^1_{\{i\}\uplus K}\big) \geq \offset\big(1^{\{i\}},0^K\big).\]
%
Now, we defined $\vec\omega^1_{\ZZ^n}$ as $\offset(1^{\{i\}})$,
so by Formula~(\ref{f4463}) it satisfies
\[ -\omega^1_{i'} + \sum_{j\neq i'} \gamma_{j-i'}(\omega_j-\omega_i') = 0 \]
for all~$i' \in \ZZ^n\setminus\{i\}$, hence \emph{a fortiori}
for all~$i'\in \ZZ^n\setminus(\{i\}\uplus K)$.
Since moreover $\vec\omega^1$ obviously coincides with~$\vec\omega^1_{\{i\}\uplus K}$
on~$\{i\} \uplus K$, this implies, by Formula~(\ref{f4463}) again, that
\[ \offset\big( \vec\omega^1_{\{i\}\uplus K} \big) =
\vec\omega^1_{\ZZ^n} = \offset\big( 1^{\{i\}} \big) .\]
So, (\ref{f3643}) becomes ``$\offset(1^{\{i\}}) \geq \offset(1^{\{i\}},0^K)$'',
what we wanted.
\end{proof}\endgroup

Thanks to Corollary~\ref{c6834} our tensorization theorems give decorrelation results
for the quadratic model:
%
\begin{Thm}\label{thmTchou}
Provided Condition~(\ref{f2380}) holds:
\begin{ienumerate}
\item\label{i4805} The quadratic model is well-$\rho$-mixing, cf.\ Definition~\ref{def-WRM}.
If $\Gamma < 1$, one can be more specific about the property ``$\rho^*(1)<1$'':
for all disjoint $I,J\subset\ZZ^n$, ${\{\vec\omega_I:\vec\omega_J\}} \leq \Gamma$.
\item\label{i4806} Moreover, if there is polynomial decay of interactions
$\gamma_{z} = O(1/|z|^{n+\alpha})$, then the model is $\alpha$-polynomially $\rho$-mixing,
and if $\gamma_{z}$ has exponential decay,
then the model is exponentially $\rho$-mixing (but not with the same rate as $\gamma_z$ in general).
\end{ienumerate}
\end{Thm}

\begin{proof}
To prove Point~(\ref{i4805}), we have to show that
\[ \sum_{z\in\ZZ^n\setminus\{0\}} \frac{a_{Q^{-1}}(z)}{a_{Q^{-1}}(0)} \leq \Gamma \]
—recall that we assumed $\Gamma < \infty$.
We write:
\[\label{f1610} \sum_{z\in\ZZ^n\setminus\{0\}} \frac{a_{Q^{-1}}(z)}{a_{Q^{-1}}(0)}
= \frac{\sum_{z\in\ZZ^n}a_{Q^{-1}}(z)}{a_{Q^{-1}}(0)} - 1 .\]
%
There, $\sum_{z\in\ZZ^n}a_{Q^{-1}}(z)$ is equal to~$1$:
indeed, $a_{Q^{-1}}$ is the convolution inverse of~$a_Q$,
so by Fubini's theorem:
\[ \sum_{z\in\ZZ^n} a_{Q^{-1}}(z) = \Big( \sum_{z\in\ZZ^n} a_Q(z) \Big)^{-1} ,\]
where
\[ \sum_{z\in\ZZ^n} a_Q(z) = \sum_{z\in\ZZ^n\setminus\{0\}}\big(-\gamma(z)\big) + (1+\Gamma)
= -\Gamma + 1 + \Gamma = 1 .\]
Now, by~(\ref{f0024}),
$a_{Q^{-1}}(0)$ is obviously bounded below by~$(1+\Gamma)^{-1}$, so in the end:
\[ \sum_{z\in\ZZ^n\setminus\{0\}} \frac{a_{Q^{-1}}(z)}{a_{Q^{-1}}(0)}
\leq \frac{1}{(1+\Gamma)^{-1}} - 1 = \Gamma .\]

To prove Point~(\ref{i4806}), we have to show that
polynomial decay of~$\gamma_z$ implies polynomial decay of~$a_{Q^{-1}}$
with the same exponent, resp.\ that exponential decay of~$\gamma_z$
implies exponential decay of~$a_{Q^{-1}}$.
This is achieved resp.\ by Lemmas~\ref{l1939a} and~\ref{l1939b} in the appendix.
\end{proof}



\section{Nonlinear lattice of particles}\label{parNonlinearLatticeOfParticles}

In this section we will consider a model with continuous spins,
but where interactions are nonlinear, so that we cannot use the properties of Gaussian variables.
One has a lattice of particles indexed by~$\ZZ^n$ (equipped with its $\ell^1$ graph structure),
each particle $i$ being described by its ``polarization'' $\omega_i\in\RR$.
Each particle is submitted to a pinning force deriving from a potential $V$,
and to interaction forces with its neighbours, the interactions deriving from a potential $W$.
In other words, the Hamiltonian is formally
\[\label{forHNL} H(\vec\omega) = \sum_{i\in\ZZ} V(\omega_i) +
\frac{1}{2} \sum_{i\sim j} W(\omega_j-\omega_i) .\]

We make the following assumptions:
\begin{Hyp}\label{hypVW}
Both $V$ and~$W$ are convex; moreover
$V$ is uniformly strictly convex and the Hessian of~$W$ is bounded,
i.e.\ there exist constants~$v_*>0$ and~$w_*<\infty$
such that for all~$x \in \RR$, $v_* \leq V''(x)$ and $W''(x) \leq w^*$.
\end{Hyp}
%
We are interested in the equilibrium state of the system at some inverse temperature
$0 < \beta < \infty$. (In the sequel we suppose that $\beta$ is fixed).

Let~$i\neq j\in\ZZ$, $K\subset\ZZ\setminus\{i,j\}$ and~$\vechat{\omega}_K \in \RR^K$;
we want to study the law of~$(\omega_i,\omega_j)$
under the law $\Pr[\Bcdot|\vec\omega_K=\vechat{\omega}_K]$.
Then, the probability distribution of the system is formally described by
\[ \dx\Pr \big( \omega_i,\omega_j,\vec\omega_{\C{K}\setminus\{i,j\}} \big) \propto
\exp \Big( -\beta
H \big(\omega_i,\omega_j,\vec\omega_{\C{K}\setminus\{i,j\}},\vechat{\omega}_K\big) \Big) .\]
%
Our assumptions ensure that the function
$H(\Bcdot,\Bcdot,\Bcdot,\vechat{\omega}_K)$ is uniformly convex,
so that the equilibrium exists and is unique.

For the sequel, we need to recall the definition of the $W_{\infty}$ Wasserstein distance:
%
\begin{Def}[see also~\cite{Winfty}]
For~$\mu_1,\mu_2$ two measures on some metric space $(X,d)$,
``$W_{\infty}(\mu_1,\mu_2) \leq \epsilon$'' means that there exists a probability measure~$\gamma$
on~$E^2$ such that the two respective marginals of~$\gamma$ are~$\mu_1$ and~$\mu_2$
and such that $d(x_1,x_2) \leq \epsilon \quad \gamma\text{-a.s.}$.
This defines a (possibly infinite) distance on the probability measures on~$E$.
\end{Def}

The fundamental lemma of this subsection is the following
\begin{Clm}\label{lem7014}
For~$\hat\omega_j \in \RR$, denote by~$\mu(\hat\omega_j)$ the law of~$\omega_i$
under~$\Pr[\Bcdot|\ab \vec{\omega}_{K\uplus\{j\}}=(\vechat{\omega}_K,\hat\omega_j)]$.
There exists a function~$\epsilon \colon \ZZ \to [0,1]$ with
$\epsilon(d) < 1$ as soon as $d > 0$ and $\epsilon(d) \stackrel{d\longto\infty}{\leq} Ce^{-\psi d}$
for some~${\psi > 0}$ and~${C < \infty}$, such that
\[ \forall \hat\omega_j^1, \hat\omega_j^2 \in \RR \qquad
W_\infty \big( \mu(\hat\omega_j^1), \mu(\hat\omega_j^2) \big)
\leq \epsilon(|j-i|) \big| \hat\omega_j^2 - \hat\omega_j^1 \big| .\]
\end{Clm}

\begin{proof}
The proof relies on the `explicit' construction of a coupling measure~$\gamma$
between~$\mu(\hat\omega_j^1)$ and~$\mu(\hat\omega_j^2)$.
To do that, we will construct $\Pr[\Bcdot|\vec{\omega}_{K\uplus\{j\}}=(\vechat{\omega}_K,\hat\omega_j)]$
thanks to a reversible Fokker--Planck dynamics, and then couple the dynamics
for~$\hat\omega_j^1$ and~$\hat\omega_j^2$.

We define the Fokker--Planck dynamics thanks to independent white noises
$(\dx B^{i'}_t)_{t\in\RR}$ for~$i' \in {\ZZ^n \setminus (K\uplus\{j\})}$.
The motion of point $i'$ is defined by:
\[ \dx\omega_{i'} = -\beta \big( V'(\omega_{i'}) +
\sum_{i''\sim i'} W'(\omega_{i'}-\omega_{i''}) \big) + \sqrt{2} \,\dx{B^{i'}_t} ,\]
with the boundary condition $\vec\omega_{K\uplus\{j\}} = \vechat{\omega}_{K\uplus\{j\}}$
for all times.
%
Coupling then consists in taking the same noise for the two processes.
The initial condition is not very important since it is asymptotically forgotten,
so we will suppose that the two systems have been coupled for an infinite time,
so that at any time both systems follow their equilibrium law.
We denote by~$\vec\omega^1(t)$ the system correponding to the boundary condition
``$\vec\omega_{K\uplus\{j\}} = (\vechat{\omega}_K,\hat\omega_j^1)$'',
resp.\ by~$\vec\omega^2(t)$ the system correponding to the other boundary condition.
We denote $\Delta_{i'}(t) \coloneqq \omega_{i'}^2(t) - \omega_{i'}^1(t)$.
Then when the dynamics are coupled,
$\vec{\Delta}$ evolves according to the following equation:
%
\[ \dx(\Delta_{i'}) =
-\beta \Big[ \big( V'(\omega^2_{i'}) - V'(\omega^1_{i'}) \big) +
\sum_{i''\sim i'} \big( W'(\omega^2_{i'}-\omega^2_{i''}) - W'(\omega^1_{i'}-\omega^1_{i''}) \big) \Big] .\]
%
Obviously the right-hand side is not a deterministic function of
$\vec{\Delta}(t)$, but it can nonetheless be written as
\[ -\beta \big[ v(i',t) \Delta_{i'}(t)
+ w(i',i'',t) \big( \Delta_{i'}(t) - \Delta_{i''}(t) \big) \big] ,\]
for some~$v(i',t)$ and~$w(i',i",t)$ satisfying
%
\begin{eqnarray}
\label{for0346v} v(i',t) & \geq & v_* \quad \text{and} \\
\label{for0346w} 0 \leq w(i',i",t) & \leq & w^* \end{eqnarray}
by Assumption~\ref{hypVW}.
%
Moreover, one has the boundary conditions:
\[ \forall t \qquad \left\{
\begin{array}{rcl} \vec{\Delta}_K &\equiv& 0 ;\\
\Delta_j &=& \hat\omega_j^2 - \hat\omega_j^1 . \end{array} \right.\]

So, $\vec{\Delta}$ is the solution of some discrete `damped heat equation',
whose coefficients can vary along time though having to satisfy bounds~(\ref{for0346v}) and~(\ref{for0346w}).
Such an equation has no stationary solution \emph{stricto sensu};
however there exists some $\vec\Delta^+$ such that
\[ \vec{\Delta}(t) \leq \vec\Delta^+ \qquad \Rightarrow \qquad
\forall t'\geq t \quad \vec{\Delta}(t') \leq \vec\Delta^+ ;\]
namely, this $\vec\Delta^+$ is defined as the solution
of the following system of equations:
$\vec\Delta^+_K \equiv 0$, $\Delta^+_j = \hat\omega_j^2 - \hat\omega_j^1$,
and for all~$i'\notin K\uplus\{j\}$,
\[\label{for7756a} 0 = - v_* \Delta^+_{i'} + \sum_{i''\sim i'}
\1{\Delta^+_{i''}\geq\Delta^+_{i'}}w^* (\Delta^+_{i''} - \Delta^+_{i'}) .\]
%
One has similarly that
\[ \vec{\Delta}(t) \geq \vec0 \qquad \Rightarrow \qquad
\forall t'\geq t \quad \vec{\Delta}(t') \geq \vec0 .\]
%
Consequently, I claim that for all~$t$ one has
\[\label{for4700} \vec0 \leq \vec\Delta(t) \leq \vec\Delta^+ :\]
indeed if the initial condition of the system satisfies~(\ref{for4700}),
then that property remains valid for all subsequent times;
now, as I told, initial conditions are asymptomatically forgotten,
so in fact (\ref{for4700}) is always satisfied.

One has the following control on~$\vec\Delta^+$:
\begin{Clm}\label{lem3945}
There exists a function~$\epsilon \colon \ZZ \to [0,1]$ with
$\epsilon(d) < 1$ as soon as $d > 0$ and $\epsilon(d) \stackrel{d\longto\infty}{\leq} Ce^{-\psi d}$
for some~$\psi > 0$ and~$C < \infty$, such that
\[ \forall i \in \ZZ^n \qquad \Delta^+_i \leq \epsilon(|j-i|) .\]
Moreover, the function~$\epsilon$ does not depend on~$K$ nor on~$j$.
\end{Clm}
%
Combining~(\ref{for4700}) with Claim~\ref{lem3945} ends the proof of Claim~\ref{lem7014}.
\end{proof}

\begingroup\def\proofname{Proof of Claim~\ref{lem3945}}
\begin{proof}
First, notice that Equation~(\ref{for7756a}) satisfies a maximum principle,
so we know in advance that $\Delta^+$ is uniquely defined with $0 \leq \Delta^+ \leq 1$ everywhere.

For~$i'\sim i''$, denote $w_{i'}(i'') \coloneqq \1{\Delta^+_{i''} \geq \Delta^+_{i'}} w^*$.
Then (\ref{for7756a}) can be rewritten into:
\[\label{for7756a+} \Delta_{i'} =
\sum_{i''\sim i'} \frac{w_{i'}(i'')}{v_*+\sum_{i''\sim i'}w_{i'}(i'')} \times \Delta_{i''}
+ \frac{v_*}{v_*+\sum_{i''\sim i'}w_{i'}(i'')} \times 0 .\]

Now I define the following Markov chain on~$\ZZ^n \uplus \{\partial\}$,
$\partial$ denoting a cemetery point:
\begin{Def}\strut
\begin{itemize}
\item If at some time the particle is on some point $i'$ of~$\ZZ^n \setminus (K\uplus\{j\})$,
at next time it jumps onto the neighbour $i''$ of~$i'$ with probability
$w_{i'}(i'') \mathbin{\big/} \big(v_*+\sum_{i''\sim i'}w_{i'}(i'')\big)$,
and it jumps onto~$\partial$ with probability
$v_* \mathbin{\big/} {\big(v_*+\sum_{i''\sim i'}w_{i'}(i'')\big)}$;
\item If the particle is somewhere in~$K \uplus \{\partial,j\}$ at some time,
then it does not move any more.
\end{itemize}
\end{Def}
%
Call~$(X_t)_{t\in\NN}$ such a Markov chain and denote by~$\mcal{L}$ its generator.
It is clear that with probability one,
$X_t$ eventually remains at some point of~$K\uplus\{\partial,i\}$.
Extend $\Delta^+$ to~$\ZZ^n\uplus\{\partial\}$ by setting $\Delta^+_\partial = 0$;
then, (\ref{for7756a+}) merely means that $\Delta^+$ is $\mcal{L}$-harmonic,
and it follows that
\[ \Delta^+_i = \EE \big[ f(X_\infty) \big| X_0=i \big] .\]
%
Thus, to bound above $\Delta^+_i$ I write that
\begin{multline} \EE\big[f(X_\infty)\big|X_0=i\big]
= \sum_{\substack{i=i_0\sim\cdots\sim i_t=j\\i_1,\ldots,i_{t-1}\notin K\uplus\{j\}}}
\,\,\prod_{u=0}^{t-1} \frac{w_{i_u}(i_{u+1})}{v_*+\sum_{i''\sim i_u}w_{i'}(i'')} \\
\leq \sum_{\substack{i=i_0\sim\cdots\sim i_t=j\\i_1,\ldots,i_{t-1}\neq j}}
\,\,\prod_{u=0}^{t-1} \frac{w_{i_u}(i_{u+1})}{\sum_{i''\sim i_u}w_{i'}(i'')}
\bigg( \frac{2dw^*}{2dw^*+v_*} \bigg)^{\!t} \\
\leq \bigg( \frac{2dw^*}{2dw^*+v_*} \bigg)^{|j-i|}\!\!
\underbrace{\sum_{\substack{i=i_0\sim\cdots\sim i_t=j\\i_1,\ldots,i_{t-1}\neq j}}
\,\,\prod_{u=0}^{t-1} \frac{w_{i_u}(i_{u+1})}{\sum_{i''\sim i_u}w_{i'}(i'')}}_{\leq 1}
\leq \bigg( \frac{2dw^*}{2dw^*+v_*} \bigg)^{\!|j-i|}.
\end{multline}
\end{proof}\endgroup

From Claim~\ref{lem7014}, we take the following
%
\begin{Cor}
For a Lipschitzian function~$f \colon \RR\to\RR$,
denote by~$\|f\|_{\mathit{Lip}}$ the optimal Lipschitz constant for~$f$.
On~$\ldb(\omega_i)$, define the (possibly infinite) norm $\|\Bcdot\|_{\mathit{Lip}}$
such that $\|f(\omega_i)\|_{\mathit{Lip}} = \|f\|_{\mathit{Lip}}$%
\footnote{This definition can be ambiguous if the support of~$\omega_i$ is not the whole $\RR$; in this case, just add an infimum in the definition.};
denote by~$\barLip(\omega_i)$ the corresponding Banach space.

Then under the law $\Pr[\Bcdot|\vec\omega_K=\vechat{\omega}_K]$,
the map $\pi_{\omega_j\omega_i}$ defined by~(\ref{for0595}) is $\epsilon(|j-i|)$-contracting
when seen as an application from $\barLip(\omega_i)$ into $\barLip(\omega_j)$.
\end{Cor}

Consequently, the map
$\pi_{\omega_i\omega_j\omega_i} \colon \barLip(\omega_i) \to \barLip(\omega_i)$
is $\epsilon(|j-i|)^2$-contracting.
But the canonical embedding $\barLip(\omega_i) \mapsto \ldb(\omega_i)$ is continuous
as our hypotheses ensure that $\Law(\omega_i)$ is uniformly log-concave,
therefore for all~$f \in \barLip(\omega_i)$ one has
\[ \limsup_{k\longto\infty}
\big|\langle \pi_{\omega_i\omega_j\omega_i}^k f, f \rangle_{\ldb(\omega_i)}\big|^{1/k}
\leq \epsilon(|j-i|)^2 .\]
%
Since $\pi_{\omega_i\omega_j\omega_i}$ is self-adjoint in~$\ldb(\omega_i)$
and $\barLip(\omega_i)$ is a dense subset of~$\ldb(\omega_i)$,
it follows by Lemma~\ref{l3253} that $\pi_{\omega_i\omega_j\omega_i}$ is
$\epsilon(|j-i|)^2$-contracting also in~$\ldb(\omega_i)$.
This, by Remark~\ref{rmk1265}, is equivalent to saying that
%
\[\label{forECRAN} \{\omega_i:\omega_j\}_{\vec\omega_K} \leq \epsilon(|j-i|) . \]

(\ref{forECRAN}) is what we need to apply Lemma~\ref{lem4067};
in the end, we get the
%
\begin{Thm}\label{thmLama}
The model~(\ref{forHNL}) is exponentially $\rho$-mixing.
\end{Thm}



\section{A hypocoercive system of interacting particles}%
\label{parAHypocoerciveSystemOfInteractingParticles}

For the time being we have only been dealing with \emph{spatial} decorrelations.
Yet I have had the idea that
the ability of Hilbertian decorrelations to get tensorized for infinite sets
could be well adapted to the study of \emph{temporal} relaxation of an infinite stochastic system:
one can consider indeed time as an extra dimension for the particle system,
which leads to a situation analogous to the parallel hyperplanes of \S~\ref{parc6177}.
In the reversible case, we saw that spectral techniques
make it possible to get $L^2$ results from $L^1$ results, cf.\ Theorem~\ref{c6177}.
Here I will show how Hilbertian decorrelations can be used
for a \emph{non-reversible} particle stochastic system.

The system which we will study here as an example
is governed by a \emph{kinetic Fokker-Planck equation}.
This equation, which arises naturally in phy\-sics, corresponds to a Hamiltonian evolution
perturbed by some noise \emph{acting on speeds}.
The study of such systems is made complicated by the fact that diffusion is only performed
along certain directions of the states space,
so that the non-reversibility of the evolution is essential to ensure convergence to equilibrium.
In~\cite{hypocoercivity}, Villani proves $L^2$ convergence for such systems
in situations where the state of the system lives in a finite-dimensional manifold.
Here we will use tensorization of Hilbertian decorrelations in a fundamental way
to get a result valid in an infinite-dimensional setting.
Moreover, we will get non-trivial bounds for arbitrary small times,
which is a new feature compared to~\cite{hypocoercivity}.

\begin{Def}\label{d7463}
For real parameters $m, \omega, c, T, \lambda > 0$%
\footnote{$m$ is the mass of each particle,
$\omega$ is the frequency corresponding to the pinning potential,
$c$ is more or less the speed of sound, expressed in inter-atomic distances by unit of time,
$T$ is the temperature and $\lambda$ is the relaxation constant of the friction.
Physical homogeneity of these constants are resp.~$[\mathsf{M}], \ab 
[\mathsf{T}^{-1}], \ab [\mathsf{T}^{-1}], \ab [\mathsf{M}\mathsf{L}^2\mathsf{T}^{-2}], \ab 
[\mathsf{T}^{-1}]$.},
we consider a system of particles $i$ indexed by~$\ZZ$,
each particle being described by its momentum $p_i \in \RR$ and its position $q_i \in \RR$.
We consider the Hamiltonian
\[ H(\vec{p},\vec{q}) = m^{-1} \sum_{i\in\ZZ} \frac{p_i^2}{2}
+ m\omega^2 \sum_{i\in\ZZ} \frac{q_i^2}{2}
+ mc^2 \sum_{i\in\ZZ} \frac{(q_{i+1}-q_i)^2}{2} .\]
Then the system $(\vec{p}(u),\vec{q}(u))$ evolves according to the Hamiltonian $H$,
plus a white noise independent on each~$p_i$, plus a friction force
$F_i = -\lambda p_i$ on each~$i$ which dissipates the energy brought by the white noise,
friction being adjusted to the noise so that their association constitutes
a (volumic) thermal bath at temperature $T$.
One computes that this means that the quadratic variation on~$p_i$ is given by
$\dx[p_i] = 2T\lambda m\,\dx{u}$.

In other words, if $\big(W_i(u)\big)_{i\in\ZZ}$ denotes a family of independent brownian motions,
the evolution of the system is given by
\[ \left\{ \begin{array}{rcl}
\dx{p_i} &=& \big( - m\omega^2 q_i + mc^2 (q_{i-1}+q_{i+1}-2q_i) - \lambda p_i \big) \,\dx{u}
+ \sqrt{2T\lambda m} \,\dx{W_i} \\
\dx{q_i} &=& m^{-1} p_i \,\dx{u} .
\end{array} \right. \]
\end{Def}

\begin{Rmk}
The system of Definition~\ref{d7463} is to be thought
as a toy model for a large class of similar systems
obtained by generalizing it in several ways.
A first example, which would change almost nothing but complicating the formalism,
is to replace the states space
$\RR\times\RR$ of each particle by~$\RR^n\times\RR^n$,
or to replace the lattice $\ZZ$ by~$\ZZ^n$.
A trickier generalization is to consider the case of non-harmonic interactions:
then I expect the results stated below to remain qualitatively true,
but proving them might be far more difficult
since one cannot use the properties of Gaussian vectors any more.
Also, if one allows for infinite-ranged interactions, which speed of decay is required
to get temporal decorrelations?

All these questions look quite worthwhile to me, though answering them
is out of the scope of this work.
Here I will only show how Hilbertian correlations make everything work fine for the toy model,
hoping that it shall be useful for the general situation.
\end{Rmk}

Let us consider the equilibrium dynamics of our system.
We fix an arbitrary time $0 < t < \infty$.
Denote by~$(p_i,q_i)$ the state of particle $i$ at time~${u=0}$,
resp.\ by~$(p'_i,q'_i)$ the state of particle $i$ at time~${u=t}$.
We have to prove the
%
\begin{Clm}\label{lem1278}
Provided $t$ is small enough, for all~$i,j\in\ZZ$ (possibly identical),
one has ${\{p_i:p'_j\}_*}, \ab {\{p_i,q'_j\}_*}, \ab {\{q_i:p'_j\}_*}, \ab 
{\{q_i,q'_j\}_*} < 1$, uniformly in~$i,j$.
Moreover, still uniformly in~$i,j$, these quantities are bounded by
$O(e^{-\gamma|j-i|})$ for some~$\gamma > 0$.
\end{Clm}

\begin{proof}
We denote by~$\eta \text{\ (resp.\ $\eta'$)} \in \RR^{\ZZ\times\{p,q\}}$
the global state $(p_i,q_i)_{i\in\ZZ}$ (resp.~$(p'_i,q'_i)_{i\in\ZZ}$) at time~$0$ (resp.~$t$).
We also denote by~$(\phi^u)_{u\geq0}$ the semigroup of operators on~$\RR^{\ZZ\times\{p,q\}}$
corresponding to the evolution of the system in absence of noise, but with the friction remaining.
Since the system is linear, the~$\phi^u$ are linear operators.

By the work of \S~\ref{parQuadraticModels},
we know that $\eta$ is distributed according to the centered Gaussian law
with covariance matrix $T^{-1}\check{C}$, where $\check{C}$ is defined as $\check{Q}^{-1}$,
the matrix $\check{Q}$ being in turn defined by:
\begin{eqnarray}
\label{forQpp} \check{Q}_{p_ip_i} &\coloneqq& m^{-1} ;\\
\label{forQpq} \check{Q}_{q_iq_i} &\coloneqq& m (\omega^2 + 2c^2) ;\\
\label{forQqq} \check{Q}_{q_iq_{i\pm1}} &\coloneqq& -mc^2 ,
\end{eqnarray}
the other entries of~$\check{Q}$ being zero.
Observe that, as the matrix of a quadratic form,
$\check{Q}$ is bounded (this is obvious from~(\ref{forQpp})--(\ref{forQqq}));
moreover, $\check{Q}^{-1}$ (actually exists and) is also bounded:
that follows from $\check{Q}$'s being bounded below by
the matrix having the same expression with~$c$ replaced by~$0$,
which we denote by~$\check{Q}^\circ$, which is a strictly positive `scalar' matrix
(modulo some homogeneity constant).

Because of the linear nature of the system, we have moreover that,
conditionally to~$\eta$, the law of~$\eta'$ is some Gaussian vector
of the form $\phi^{t}\eta + \theta$, where $\theta$ is a centered Gaussian vector
whose law does not depend on~$\eta$.
Let us denote by~$\hat{C}$ the covariance matrix of~$\theta$,
and $\hat{Q} = \hat{C}^{-1}$ —though for the time being it is not clear that $\hat{Q}$ exists.

Then, we can formally write the covariance matrix $\bar{C}$ of~$(\eta,\eta')$ as
$\bar{C} = \bar{Q}^{-1}$, with:
\[\label{for7743}
\bar{Q}(\eta,\eta') = \check{Q}(\eta) + \hat{Q}(\eta'-\phi^{t}\eta) .\]
(Note that $\bar{Q}$ is a quadratic form on~$\RR^{\ZZ\times\{p,q,p',q'\}}$,
while $\check{Q}$ and~$\hat{Q}$ were defined on~$\RR^{\ZZ\times\{p,q\}}$).

\begin{Not}
In the sequel, we shorthand ``$\ZZ\times\{p,q\}$'' into~``$\ZZ^{\uplus2}$'',
resp.\ ``$\ZZ\times\{p,q,p',q'\}$'' into~``$\ZZ^{\uplus4}$''.
\end{Not}

Now I claim that there exists
constants $0<r\leq R<\infty$ such that
$r\mathbf{I} \leq \bar{Q} \leq R\mathbf{I}$.
Well, this is meaningless \emph{stricto sensu}, because all the entries of~$\bar{Q}$
do not have the same physical homogeneity,
so we have to `convert' momenta into positions by dividing them
by some homogeneity parameter~$\chi$, say $\chi = m\omega$
—but other choices may be more relevant.

First, I claim that $\bar{Q} \geq \frac{1}{2}(\chi^2 m^{-1}\wedge m\omega^2) \mbf{I}$.
Let indeed $(\eta,\eta') = (\vec{p}_{\ZZ},\vec{q}_{\ZZ},\vec{p'}_{\ZZ},\vec{q'}_{\ZZ})
\in \RR^{\ZZ^{\uplus4}}$ with finite support.
We observe that
\[ \|(\eta,\eta')\|^2 = \sum_{i\in\ZZ} \big(\chi^{-2}p_i^2+q_i^2+\chi^{-2}{p'}_i^2+{q'}_i^2\big)
= \| \eta \|^2 + \| \eta' \|^2 ,\]
so that either $\|\eta\|^2 \geq \frac{1}{2} \|(\eta,\eta')\|^2$
or $\|\eta'\|^2 \geq \frac{1}{2} \|(\eta,\eta')\|^2$.
Now, recalling the definition of~$\check{Q}^\circ$ a few lines above,
$\check{Q}(\eta) \geq \check{Q}^\circ(\eta) \geq (\chi^2 m^{-1}\wedge m\omega^2) \|\eta\|^2$,
so by~(\ref{for7743}), $\bar{Q}(\eta,\eta') \geq (\chi^2 m^{-1}\wedge m\omega^2) \|\eta\|^2$.
Since reversing the direction of time yields the same system with the sign of speeds reversed, which does not change the norms of~$\eta$ and~$\eta'$,
one has similarly $\bar{Q}(\eta,\eta') \geq (\chi^2 m^{-1}\wedge m\omega^2) \|\eta'\|^2$.
The claim follows.

The second point consists in proving that $\bar{Q}$ is bounded above.
On the one hand, by~(\ref{forQpp})--(\ref{forQqq}),
\[ \check{Q}(\eta) \leq \big(m^{-1}\chi^2\vee m(\omega^2+4c^2)\big) \|\eta\|^2
\leq \big(m^{-1}\chi^2\vee m(\omega^2+4c^2)\big) \|(\eta,\eta')\|^2 .\]

Next, the difficult point is to prove that $\hat{Q}(\eta'-\phi^{t}\eta)$
(exists and) can be bounded above by a multiple of~$\|(\eta,\eta')\|^2$.
We begin with transforming the original problem
of bounding a quadratic form on~$\RR^{\ZZ^{\uplus4}}$
into a problem on~$\RR^{\ZZ^{\uplus2}}$.
Indeed, $\|\phi^t\eta\|$ is bounded by a multiple of~$\|\eta\|$,
since the operator $\phi^t$ dissipates the energy $H(\eta)$,
energy which the previous work on~$\check Q$
proved to be controlled below and above by~$\|\eta\|^2$;
therefore, it suffices to prove that the quadratic form $\hat{Q}(\eta)$ on~$\RR^{\ZZ^{\uplus2}}$
is bounded by a multiple of~$\|\eta\|^2$ to achieve our goal.

The natural quantity to be computed for~$\theta$
(recall that $\theta$ denotes the total effect of noise between times~$0$ and~$t$)
is its covariance matrix $\hat{C}$.
Its expression is the following (the notation is explained just below):
\[\label{for5483} \hat{C} = 2T\lambda m \int_0^{t} \phi^{t-u} \mbf{I}_p \T{(\phi^{t-u})} \,\dx{u} ,\]
where $\mbf{I}_p$ is the diagonal matrix being~$1$ on diagonal entries indexed by some~$p_i$
and~$0$ on diagonal entries indexed by some~$q_i$,
and $\T{(\phi^{t-u})}$ is the transpose of the linear operator $\phi^{t-u}$
seen as a square matrix indexed by~$\ZZ^{\uplus2}$.
This decomposition means that we are summing the contributions
of all the elementary noises occuring at times $u \in [0,t]$,
using that these elementary noises are independent.

Now we need an approximate expression for~$\phi^u$, $u\in[0,t]$.
Here for the sake of legibility I will remain at a formal level, giving only limited expansions;
it is essential nevertheless to keep in mind that all the ``$O(*)$''
can be made explicit by using Gronwall's lemma,
and that these explicit values ensure that the~$O(*)$ behave well provided $t$ is small enough.
%
One finds that
\begin{eqnarray}
\label{for5476a}
\phi^u \delta_{p_i} \cdot p_j &=& c^{2|j-i|} \frac{u^{2|j-i|}}{(2|j-i|)!} + O(u^{2|j-i|+2}) ;\\
\phi^u \delta_{p_i} \cdot q_j &=& m^{-1}c^{2|j-i|} \frac{u^{2|j-i|+1}}{(2|j-i|+1)!}
+ O(u^{2|j-i|+3}) ;\\
\phi^u \delta_{q_i} \cdot p_i &=& m\omega^2 u + O(u^3) ;\\
\phi^u \delta_{q_i} \cdot p_{j\neq i} &=& mc^{2|j-i|} \frac{u^{2|j-i|-1}}{(2|j-i|-1)!}
+ O(u^{2|j-i|+1}) ;\\
\label{for5476z}
\phi^u \delta_{q_i} \cdot q_j &=& c^{2|j-i|} \frac{u^{2|j-i|}}{(2|j-i|)!} + O(u^{2|j-i|+2}) .
\end{eqnarray}
%
Injecting Equations~(\ref{for5476a})--(\ref{for5476z}) into~(\ref{for5483}), one finds that:%
\footnote{Recall that $\hat{C}$, as a covariance matrix, is symmetric.}
\begin{eqnarray}
\label{for5579a}
\hat{C}_{p_ip_i} &=& 2T\lambda m t + O(t^3) ;\\
\hat{C}_{p_iq_i} &=& T\lambda t^2 + O(t^4) ;\\
\hat{C}_{q_iq_i} &=& \mathsmaller{\frac{2}{3}}T\lambda m^{-1} t^3 + O(t^5) ;\\
\hat{C}_{p_ip_{j\neq i}} &=& O(t^{2|j-i|+1}) ;\\
\hat{C}_{p_iq_{j\neq i}} &=& O(t^{2|j-i|+2}) ;\\
\label{for5579z}
\hat{C}_{q_iq_{j\neq i}} &=& O(t^{2|j-i|+3}).
\end{eqnarray}

Consequently, the covariance matrix $\hat{C}$ can be seen
as a perturbation of the matrix $\hat{C}^\circ$
which is defined by Equations~(\ref{for5579a})--(\ref{for5579z}),
but with the ``$O(*)$'' terms replaced by~$0$.
%
Since $\hat{C}^\circ$ is invertible, with an explicitly computable inverse,
one finds that $\hat{C}$ is invertible too with:
%
\begin{eqnarray}
\hat{Q}_{p_ip_i} &=& 2T^{-1}\lambda^{-1}m^{-1}t^{-1} + O(t) ;\\
\hat{Q}_{p_iq_i} &=& -3T^{-1}\lambda^{-1}t^{-2} + O(1) ;\\
\hat{Q}_{q_iq_i} &=& 6T^{-1}\lambda^{-1}mt^{-3} + O(t^{-1}) ;\\
\hat{Q}_{p_ip_{j\neq i}} &=& O(t^{2|j-i|-1}) ;\\
\hat{Q}_{p_iq_{j\neq i}} &=& O(t^{2|j-i|-2}) ;\\
\hat{Q}_{q_iq_{j\neq i}} &=& O(t^{2|j-i|-3}).
\end{eqnarray}
%
In the end, provided that $t$ is small enough, we have proved that
$\hat{Q}(\eta)/\|\eta\|^2 \leq 6T^{-1}\lambda mt^{-3} + O(t^{-1}) \ab {< \infty}$.

Actually we have proved more than that:
not only we have a bound on the operator norm of~$\bar{Q}$,
but we have bounded it entry-wise.
More precisely, expanding the~$O(*)$, we find that provided $t$ is small enough,
there exists constants~$A < \infty$ and~$\gamma > 0$ such that for all~$i,j\in\ZZ$,
\[ \underbrace{\hat{Q}_{p_ip_j}, \hat{Q}_{p_iq_j}, \ldots, \hat{Q}_{q'_iq'_j}}%
_{\text{all 16 possibilities}} \leq A e^{-\gamma|j-i|} .\]

\begin{Not}
From now on we denote the basic variables $p_i, q_i, p'_i, q'_i$ of our system by~$X_i$,
$i \in \ZZ^{\uplus4}$.
\end{Not}

Now the question is: for~$i\neq j \in \ZZ^{\uplus4}$,
$K \subset \ZZ^{\uplus4} \setminus \{i,j\}$,
what is the value of~${\{X_i:X_j\}}_{\vec{X}_K}$?
%
By the properties of Gaussian variables [Theorem~\ref{pro1857}], the answer is the following.
Let~$\bar{Q}_{|_{\ZZ^{\uplus4}\setminus K}}$ be the restriction
of~$\bar{Q}$ to indexes in~$(\ZZ^{\uplus4}\setminus K)$.
Since $r\mathbf{I} \leq \bar{Q} \leq R\mathbf{I}$, the same holds
for~$\bar{Q}_{|_{\ZZ^{\uplus4}\setminus K}}$, so this matrix is invertible;
denote by~$\bar{C}^{|^{\ZZ^{\uplus4}\setminus K}}$ its inverse.
This matrix is the covariance matrix of (the centered version of)
$\vec{X}_{\ZZ^{\uplus4}\setminus K}$ under some fixed value for~$\vec{X}_K$;
thus:
\[\label{for4972} \{X_i:X_j\}_{\vec{X}_K} =
\frac{\Big| \bar{C}^{|^{\ZZ^{\uplus4}\setminus K}}_{ij} \Big|}
{\sqrt{\bar{C}^{|^{\ZZ^{\uplus4}\setminus K}}_{ii}
\bar{C}^{|^{\ZZ^{\uplus4}\setminus K}}_{jj}}} .\]

It remains to control the entries of~$\bar{C}^{|^{\ZZ^{\uplus4}\setminus K}}$, uniformly in~$K$.
We need two types of control: first an exponential control when $i$ is far away from $j$,
then a non-trivial control for the values of~$i$ and~$j$ corresponding to close
(or even identical) atoms.

Let us start with the first one.
$\bar{C}^{|^{\ZZ^{\uplus4}\setminus K}}_{ii}$ and
$\bar{C}^{|^{\ZZ^{\uplus4}\setminus K}}_{jj}$ are bounded below by~$R^{-1}$,
so we just have to bound above $\bar{C}^{|^{\ZZ^{\uplus4}\setminus K}}_{ij}$.
This is achieved by a direct use of Lemma~\ref{lem4060} in appendix.

Concerning the uniform non-trivial control, since $r\mbf{I} \leq \bar{Q} \leq R\mbf{I}$
one has $r\mbf{I} \leq \bar{Q}_{|_{\ZZ^{\uplus4}\setminus K}} \leq R\mbf{I}$,
hence $R^{-1}\mbf{I} \leq \bar{C}^{|^{\ZZ^{\uplus4}\setminus K}} \leq r^{-1}\mbf{I}$,
hence $R^{-1}\mbf{I} \leq \big(\bar{C}^{|^{\ZZ^{\uplus4}\setminus K}}\big)_{|_{\{i,j\}^2}} \leq r^{-1}\mbf{I}$;
from this and~(\ref{for4972}),
\[ \forall i \neq j \in I \qquad \{X_i:X_j\}_* \leq \frac{R-r}{R+r} < 1 .\]
%
\end{proof}

From Claim~\ref{lem1278}, we get the main result of this subsection:
%
\begin{Thm}\label{thmParis}
For the model of Definition~\ref{d7463}, for all~$t>0$, $\{\eta,\eta'\} < 1$.
\end{Thm}

\begin{proof}
First, if $t$ is small enough so that Claim~\ref{lem1278} holds,
direct application of Lemma~\ref{lem4067} proves the result,
as the~$\{X_i,X_j\}_*$ are summable (since they decrease exponentially)
and they all are $<1$.

Now for larger $t$, fix some $0<t_1<t$ so that Claim~\ref{lem1278} holds for~$t_1$.
Then we notice that $\eta \to \eta(t_1) \to \eta'$ is a Markov chain
(with ``$\eta(t_1)$'' standing for ``$\big(\vec{p}(t_1),\vec{q}(t_1)\big)$''),
so by Proposition~\ref{cor0706}, $\{\eta,\eta'\} \leq \{\eta,\eta(t_1)\} < 1$.
\end{proof}




\section{Appendix: Inverses of `nearly diagonal' matrices}%
\label{parInversesOfDiagonalConcentratedMatrices}

The goal of this appendix is to state and prove a few lemmas sharing the same spirit:
``\emph{if a matrix is `nearly diagonal', then it shall be invertible
and its inverse shall also be `nearly diagonal' with the same type of decay}''.

\subsection{Matrices with exponential decay}

The goal of this subsection is to prove the following
%
\begin{Lem}\label{lem4060}
Let~$I \subset \ZZ$ and let~$(\!(M_{ij})\!)_{(i,j)\in I^2}$ be a matrix.
Assume that, when seen as a quadratic form on~$L^2(I)$, one has
$r\mbf{I} \leq M \leq R\mbf{I}$ for $0<r\leq R<\infty$ —in particular, $M$ is invertible.
Assume moreover that there exists constants~$A < \infty$ and~$\gamma > 0$ such that
for all~$i,j\in I$, $|M_{ij}| \leq A e^{-\gamma|j-i|}$.

Then there exist constants~$A' < \infty$ and~$\gamma' > 0$
which are explicit functions of~$r,R,\gamma,A$ (so \emph{they do not depend on~$I$}),
such that one has the following control on the entries of~$M^{-1}$:
\[ \forall i,j \in I \qquad (M^{-1})_{ij} \leq A' e^{-\gamma'|j-i|} .\]
\end{Lem}

\begin{proof}
Up to multiplying by a scalar, one can assume that $R=1$.
Then $M$ writes $M = \mbf{I} - H$, where $0 \leq H \leq (1-r) \mbf{I}$;
since $H$ is symmetric, that inequality means that $\VERT H \VERT \leq 1-r < 1$.
Therefore, for all~$k \in \NN$ one has $\VERT H^k \VERT \leq (1-r)^k$,
which allows us to write $M^{-1}$ as a series expansion:
\[ M^{-1} = \sum_{k=0}^\infty H^k .\]

Up to replacing $A$ by~$A+1$,
we have the same entry-wise control on~$H$ as on~$M$.
Then one sees by induction that for all~$k\in\NN$,
\[\label{for5051}
\forall i,j\in I \qquad \big| (H^k)_{ij} \big| \leq A_1^k e^{-\gamma_1|j-i|} ,\]
where $\gamma_1$ is an arbitrary parameter in~$(0,\gamma)$ and
\[ A_1 \coloneqq \sum_{z\in\ZZ} Ae^{-\gamma|z| + \gamma_1 z}
= \frac{(1-e^{-2\gamma})A}{(1-e^{-(\gamma-\gamma_1)})(1-e^{-(\gamma+\gamma_1)})} \]
—observe that $I$ does not appear in the expression of~$A_1$.
%
Since $A_1$ is greater than $1$,
(\ref{for5051}) is not enough to get an entry-wise control on~$M^{-1}$.
But now observe that the bound $\VERT H^k \VERT \leq (1-r)^k$ implies
that all the~$(H^k)_{ij}$ are bounded by~$(1-r)^k$ in absolute value;
thus:
\[ \big| (M^{-1})_{ij} \big| \leq
\sum_{k=0}^\infty \big( e^{-\gamma_1|j-i|} A_1^k \wedge (1-r)^k \big) \\
\leq \bigg( \frac{A_1}{A_1-1} + \frac{1}{r} \bigg)
\exp \left( -\frac{|\ln(1-r)| \gamma_1}{|\ln(1-r)| + \ln A_1} |j-i| \right) ,
\]
from which you read suitable values for~$A'$ and~$\gamma'$.
\end{proof}


\subsection{Convolution inverses of rapidly decreasing functions}%
\label{parDeconvolutionOfRapidlyDecreasingFunctions}

\begin{NOTA}
In all this subsection, we work on~$\ZZ^n$ for some $n\in\NN^*$;
$\RR^n$ is endowed with some fixed norm~$|\Bcdot|$.
\end{NOTA}

\begin{Rmk}
Here I will deal with fonctions on~$\ZZ^n$, but the results of this subsection
could also be tranposed for functions on~$\RR^n$.
\end{Rmk}

\begin{Def}
If $a:\ZZ^n\to\RR$ is some integrable function with
$\|a\|_{\ell^1} < 1$, we define
\[\label{f8928} B[a] = a + a* a + a* a* a + \cdots ,\]
which is the sum of a convergent series in~$\ell^1(\ZZ^n)$.
$B[a]$ is the function~$b\in\ell^1(\ZZ^n)$ characterized by:
\[ (\delta_0-a) * (\delta_0+b) = \delta_0 .\]
\end{Def}

\begin{Def}\label{def726}
A function~$a:\ZZ^n\to\RR$ is said to have \emph{exponential decay} if there exists some $\beta > 0$
such that, for all~$\beta'<\beta$, $a(z) = O(e^{-\beta'|z|})$ when~$|z|\longto\infty$.
The minimal $\beta$ satisfying that property is called the
\emph{(exponential) rate of decay} of~$a$.
\end{Def}

\begin{Lem}\label{l1939a}
Let~$a\in\ell^1(\ZZ^n)$ with $\|a\|_{\ell^1} < 1$.
If $a$ has exponential decay, then so does $B[a]$.
\end{Lem}

\begin{proof}
Denoting by~$|a|$ the function defined by~$|a|(z) = |a(z)|$,
it is clear by~(\ref{f8928}) that
\[ \forall z\in\ZZ^n \qquad \big|B[a](z)\big| \leq B\big[|a|\big](z) ,\]
therefore it suffices to prove the case where $a$ is nonnegative.
In that case, $B[a]$ will also be nonnegative.

Let~$(\RR^n)^*$ denote the dual space of~$\RR^n$, endowed with the dual norm
\[ \forall \lambda \in (\RR^n)^* \qquad
|\lambda|_* = \sup_{\substack{z\in\RR^n\\|z|=1}} |\langle \lambda, z \rangle| .\]
For a nonnegative function~$a$, we define its Laplace transform
$\mcal{L}\{a\} \colon (\RR^n)^* \to \RR_+ \cup \{+\infty\}$ by
\[\label{f9677} \mcal{L}\{a\}(\lambda) = \sum_{z\in\ZZ^n} e^{\langle\lambda,z\rangle} a(z) .\]
Then, saying that $a$ has exponential decay with rate $\gamma$
is equivalent to saying that, for all~$\lambda \in (\RR^n)^*$ with $|\lambda|_* < \gamma$,
$\mcal{L}\{a\}(\lambda)$ is finite.

Since Laplace transform is linear and turns convolution into ordinary product,
(\ref{f8928}) yields, for all~$\lambda \in (\RR^n)^*$:
\[\label{f9574} \mcal{L}\{B[a]\}(\lambda) = \mcal{L}\{a\}(\lambda) + \mcal{L}\{a\}(\lambda)^2
+ \mcal{L}\{a\}(\lambda)^3 + \cdots, \]
which converges if and only if $\mcal{L}\{a\}(\lambda) < 1$.

Now, since $a$ is nonnegative, by~(\ref{f9677}) the function~$\mcal{L}\{a\}$ is convex,
so it is continuous on the interior of the domain where it is finite.
By the exponential decay hypothesis, that domain contains a neighbourhood of~$0$,
so $\mcal{L}\{a\}$ is continuous at~$0$.
And since $\mcal{L}\{a\}(0) = \sum_{z\in\ZZ^n} a(z) = \|a\|_{\ell^1} < 1$,
there is a neighbourhood of~$0$ on which $\mcal{L}\{a\}<1$ and thus $\mcal{L}\{B[a]\}<\infty$.
This implies that $B[a]$ has exponential decay.
\end{proof}

\begin{Rmk}\label{rmk9890}
This proof also shows that (for nonnegative $a$)
the rate of decay of~$B[a]$ will never be greater than the rate of decay of~$a$.
In general, it is even strictly smaller,
since all the values of~$\lambda$ for which $1\leq\mcal{L}\{a\}(\lambda)<\infty$ yield
a finite Laplace tranform for~$a$ but an infinite one for~$B[a]$.
For example, take $n=1$ and $a = e^{-1}\delta_{1}$,
which has exponential decay with infinite rate since it is compactly supported;
then the $k$-th convolution power of~$a$ is $a^{*k} = e^{-k}\delta_k$,
so that $B[a]$ is the function
\[ B[a](z) = \1{z>0} e^{-z} ,\]
which also has exponential decay, but with rate $1$ only.
\end{Rmk}

\begin{Lem}\label{l1939b}
If $\|a\|_{\ell^1(\ZZ^n)} < 1$ and $a(z) = O(1/|z|^\alpha)$ when~$|z|\longto\infty$
for some $\alpha > n$, then $B[a](z) \ab {= O(1/|z|^\alpha)}$ when~$|z|\longto\infty$.
\end{Lem}

\begin{proof}
Let~$a$ satisfy the assumptions of the lemma for some $\alpha$.
Like in the proof of Lemma~\ref{l1939b}, we can assume that $a$ is nonnegative.
For~$d > 0$, we define the function~$\phi_d \colon \ZZ^n\to\RR$ by:
\[ \phi_d(z) \coloneqq 1/(|z|\wedge d)^\alpha ,\]
which is in~$\ell^1(\ZZ^n)$ since $\alpha > n$.
Then the key claim is the following sub-lemma, whose proof is postponed:
%
\begin{Lem}\label{c9125}
Under the assumptions of Lemma~\ref{l1939b},
there exists some $\rho < 1$ and some $d \in (0,\infty)$ such that, pointwise,
\[\label{f9148} \phi_d * a \leq \rho\phi_d .\]
\end{Lem}

Admitting Lemma~\ref{c9125}, take~$\rho$ and~$d$ such that (\ref{f9148}) is satisfied.
The assumption on~$a$ implies that there exists some $C < \infty$ such that $a \leq C\phi_d$;
therefore by~(\ref{f9148}) one also has $a * a \leq C \phi_d * a \leq \rho C \phi_d$,
whence by~(\ref{f9148}) again $a * a * a \leq \rho C \phi_d * a \leq \rho^2 C \phi_d$,
etc.. In the end,
\[ B[a] \leq C \phi_d + \rho C \phi_d + \rho^2 C \phi_d + \cdots \leq \frac{C}{1-\rho} \phi_d ,\]
which implies that $B[a](z) = O(1/|z|^\alpha)$.
\end{proof}

\begingroup\def\proofname{Proof of Lemma~\ref{c9125}}
\begin{proof}
Denote $S \coloneqq \|a\|_{\ell^1}$, which by hypothesis is $<1$,
and fix $\epsilon \in (0,1/2)$ such that $(1-\epsilon)^\alpha > S$.
Let~$d \in (0,\infty)$, devised to be quite large; our goal is to bound above
$\big(\phi_d * a\big)(z)$ for all~$z \in \ZZ^n$.
Since $\phi_d$ is bounded above by~$d^{-\alpha}$, one has obviously for all~$z \in \ZZ^n$:
\[\label{f2348} \big(\phi_d * a\big)(z) \leq d^{-\alpha} \sum_{z\in\ZZ^n} a(z) = S d^{-\alpha} ,\]
whence $\big(\phi_d * a\big)(z) \leq S \phi_d(z)$ for all~$z$ with $|z|\leq d$.
Since $S<1$, the claim is therefore okay for $|z|\leq d$. 

Now, let~$z\in\ZZ^n$ with $|z|>d$. We have to bound above
\[ \big(\phi_d*a\big)(z) = \sum_{\substack{x,y\in\ZZ^n\\x+y=z}} \phi_d(x)\,a(y) .\]
We decompose this sum into three pieces:
\[ \big(\phi_d*a\big)(z) = \sum_{|y|\leq\epsilon|z|} \phi_d(z-y)\,a(y) +
\sum_{\substack{|x|>\epsilon|z|\\|z-x|>\epsilon|z|}} \phi_d(x)\,a(z-x) +
\sum_{|x|\leq\epsilon|z|} \phi_d(x)\,a(z-x) ,\]
which we shorthand into ``$\onecirc + \twocirc + \threecirc$''.

We bound these three terms separately.
For~\onecirc, we observe that for $|y| \leq \epsilon|z|$, $|z-y| \geq (1-\epsilon)|z|$
by the triangle inequality, thus $\phi_d(z-y) \leq {\big((1-\epsilon)|z|\big)^{-\alpha}} \ab 
= (1-\epsilon)^{-\alpha} \phi_d(z)$, whence by summing:
\[\label{f1034a} \onecirc \leq (1-\epsilon)^{-\alpha} \phi_d(z) \sum_{|y|\leq\epsilon|z|} a(y)
\leq (1-\epsilon)^{-\alpha} S \phi_d(z) .\]
Similarly, for $|x| \leq \epsilon|z|$, $C$ denoting a constant such that $a\leq C\phi_d$,
one has ${a(z-x)} \ab \leq C\big((1-\epsilon)|z|\big)^{-\alpha}$, thus:
\[\label{f1034b} \threecirc \leq (1-\epsilon)^{-\alpha} C \|\phi_d\|_{\ell^1} \, \phi_d(z) .\]
Of course, $\|\phi_d\|_{\ell^1}$ depends on~$d$; the important point is that,
by dominated convergence, $\|\phi_d\|_{\ell^1} \longto 0$ when~$d \longto \infty$.

Finally, provided $d$ is large enough, Term~\twocirc\ will be well approximated by an integral:
\[\label{f1820} \twocirc \leq
\sum_{\substack{x\in\ZZ^n\\|x|,|z-x|>\epsilon|z|}} \frac{1}{|x|^\alpha}\times\frac{C}{|z-x|^\alpha}
\simeq \int_{\begin{subarray}{l}x\in\RR^n\\|x|,|z-x|>\epsilon|z|\end{subarray}}
\frac{C}{(|x|\,|z-x|)^\alpha}\, \dx{x} ,\]
where ``$\simeq$'' means that the ratio between the quantites at each side of that symbol
can be made arbitrarily close to~$1$ when~$d \longto \infty$, uniformly in~$z$.
%
Indeed, the difference between the sum and the integral
is due to two causes: first, approximating the integral on a unit square of~$\RR^n$
by the value of the integrand at the center of this square, second, summing (or not summing)
terms of the discrete sum corresponding to squares
that are not entirely in the domain of the integral.
%
For the first cause,
on the domain of the integral, $C/(|x||z-x|)^\alpha$
varies of at most $O(1/|z|)$ in relative value on all the unit squares.
For the second cause, the border of the domain of the integral
is made of two $(n-1)$-dimensional spheres of radius $\epsilon |z|$,
so it crosses $O(|z|^{n-1})$ unit squares. Since $C/(|x||z-x|)^\alpha$
is bounded by~$C(\epsilon(1-\epsilon))^{-\alpha} |z|^{-2\alpha}$
on the domain of the integral, the (absolute) error due to boundary squares
is at most $O(|z|^{n-1-2\alpha})$. As the integral itself is proportional to
$|z|^{n-2\alpha}$ (cf.\ the change of variables below),
the relative error due to boundary squares is at most $O(1/|z|)$ too,
and $O(1/|z|) = o(1)$ since $|z| > \epsilon d$.

Making the change of variables $x = |z|x'$, (\ref{f1820}) becomes:
\[ \twocirc \lesssim C|z|^{n-2\alpha}
\int_{|x'|,|1-x'|>\epsilon} \frac{1}{|x'|^\alpha|1-x'|^\alpha} \, \dx{x'} ,\]
which I shorthand into ``$\twocirc \lesssim \mcal{I}C|z|^{n-2\alpha}$''.
Since $|z|>d$ and $\alpha > n$, this bound implies:
\[\label{f1034c} \twocirc \lesssim \frac{\mcal{I}C}{d^{\alpha-n}} \, \phi_d(z) .\]

Combining~(\ref{f1034a}), (\ref{f1034b}) and~(\ref{f1034c}),
one finally gets that when~$d\longto\infty$, for all~$|z|\geq d$,
\[ \big(\phi_d*a\big)(z) \leq \rho(d) \, \phi_d(z) ,\]
with
\[ \rho(d) = (1-\epsilon)^{-\alpha} (S+C\|a\|_{\ell^1})
+ (1+o(1)) \mcal{I}C/d^{\alpha-n} .\]
$\rho(d)$ tends to~$(1-\epsilon)^{-\alpha} S < 1$ when~$d\longto\infty$,
so it is $<1$ provided $d$ is large enough, which is what we wanted.\linebreak[1]\strut
\end{proof}\endgroup

